\documentclass[../../main.tex]{subfiles}% 注意这里的文件路径不能用 ./main.tex ,否则用latexmk编译子文件会报错
\graphicspath{{\subfix{./image/}}{\subfix{../../image/}}} % 指定图片引用路径,后续可以直接使用图片文件名
% 如果图片存放在上级目录,则使用 ../../image/ 作为备用路径

% 例如:
% \begin{figure}[H]
% \centering
% \includegraphics[width=0.3\textwidth]{图.png}
% \caption{}
% \label{figure:图}
% \end{figure}
% 注意:上述\label{}一定要放在\caption{}之后,否则引用图片序号会只会显示??.

\begin{document}

\section{定积分}

\subsection{建立积分递推}

\begin{example}
计算\(\int_{0}^{\frac{\pi}{2}} \cos^{n}x\sin(nx)\mathrm{d}x\)，\(n \in \mathbb{N}\)。 
\end{example}
\begin{proof}
利用分部积分和和差化积公式可得
\begin{align*}
I_n&=\int_0^{\frac{\pi}{2}}\cos ^nx\sin(nx) \mathrm{d}x=\int_0^{\frac{\pi}{2}}\cos ^{n - 1}x\cdot\cos x\sin(nx) \mathrm{d}x\\
&=\frac{1}{2}\int_0^{\frac{\pi}{2}}\cos ^{n - 1}x\cdot[\sin((n + 1)x)+\sin((n - 1)x)] \mathrm{d}x\\
&=\frac{I_{n - 1}}{2}+\frac{1}{2}\int_0^{\frac{\pi}{2}}\cos ^{n - 1}x[\sin(nx)\cos x+\cos(nx)\sin x] \mathrm{d}x\\
&=\frac{I_{n - 1}}{2}+\frac{I_n}{2}-\frac{1}{2}\int_0^{\frac{\pi}{2}}\cos ^{n - 1}x\cos(nx) \mathrm{d}\cos x\\
&=\frac{I_{n - 1}+I_n}{2}-\frac{1}{2n}\int_0^{\frac{\pi}{2}}\cos(nx) \mathrm{d}\cos ^nx\\
&=\frac{I_{n - 1}+I_n}{2}+\frac{1}{2n}-\frac{1}{2}\int_0^{\frac{\pi}{2}}\cos ^nx\sin(nx) \mathrm{d}x\\
&=\frac{I_{n - 1}+I_n}{2}+\frac{1}{2n}-\frac{I_n}{2}\\
&=\frac{I_{n - 1}}{2}+\frac{1}{2n}.
\end{align*}
故\(I_n=\frac{I_{n - 1}}{2}+\frac{1}{2n}\)，则两边同乘\(2^n\)（强行裂项）
\begin{align*}
2^nI_n = 2^{n - 1}I_{n - 1}+\frac{2^{n - 1}}{n},n = 1,2,\cdots
\end{align*}
又注意到\(I_0 = 0\)，从而
\begin{align*}
2^nI_n = 0+\sum_{k = 1}^n\frac{2^{k - 1}}{k}\Rightarrow I_n=\frac{1}{2^n}\sum_{k = 1}^n\frac{2^{k - 1}}{k}.
\end{align*}
\end{proof}

\begin{proposition}\label{example:经典定积分(必记)}
设$n\in \mathbb{N}$,证明:
\begin{enumerate}[(1)]
\item $\int_0^{\pi}{\frac{\sin \left( nx \right)}{\sin x}\mathrm{d}x}= 
\begin{cases}
0, & n\text{为偶数} \\
\pi, & n\text{为奇数}
\end{cases}.$

\item $\int_0^{\pi}{\frac{\sin ^2\left( nx \right)}{\sin ^2x}}\mathrm{d}x=n\pi$.

\item $\int_0^{\pi}{\frac{\sin ^2\left( nx \right)}{\sin x}}\mathrm{d}x=\sum_{k=1}^n{\frac{2}{2k-1}}$.
\end{enumerate}
\end{proposition}
\begin{note}
提示:$\sin ^2x-\sin ^2y=\sin \left( x-y \right) \sin \left( x+y \right)$(证明见\refpro{theorem:三角平方差}).
\end{note}
\begin{proof}
\begin{enumerate}[(1)]
\item 记\(I_n = \int_0^{\pi} \frac{\sin(nx)}{\sin x} \, \mathrm{d}x\)，则
\begin{align*}
I_{n+2} - I_n = \int_0^{\pi} \frac{\sin((n+2)x) - \sin(nx)}{\sin x} \, \mathrm{d}x 
= \int_0^{\pi} \frac{2\cos((n+1)x) \sin x}{\sin x} \, \mathrm{d}x 
= 2\int_0^{\pi} \cos((n+1)x) \, \mathrm{d}x 
= 0.
\end{align*}
于是
\begin{align*}
\int_0^{\pi} \frac{\sin(nx)}{\sin x} \, \mathrm{d}x = I_n = I_{n-2} = \cdots = 
\begin{cases}
I_0, & n\text{为偶数} \\
I_1, & n\text{为奇数}
\end{cases} = 
\begin{cases}
0, & n\text{为偶数} \\
\pi, & n\text{为奇数}
\end{cases}.
\end{align*}

\item 记 \( I_n = \int_0^\pi \frac{\sin^2(nx)}{\sin^2 x} \, \mathrm{d}x \)，则
\begin{align}
I_{n+1} - I_n &= \int_0^\pi \frac{\sin^2((n+1)x) - \sin^2(nx)}{\sin^2 x} \, \mathrm{d}x = \int_0^\pi \frac{\sin x \cdot \sin((2n+1)x)}{\sin^2 x} \, \mathrm{d}x \nonumber \\
&= \int_0^\pi \frac{\sin((2n+1)x)}{\sin x} \, \mathrm{d}x \xlongequal{\text{\nrefpro{example:经典定积分(必记)}{(1)}}} \pi.
\end{align}
于是
\[
\int_0^\pi \frac{\sin^2(nx)}{\sin^2 x} \, \mathrm{d}x = I_n = \pi + I_{n-1} = \cdots = (n-1)\pi + I_1 = n\pi.
\]

\item 记 \( I_n = \int_0^\pi \frac{\sin^2(nx)}{\sin x} \, \mathrm{d}x \)，则
\begin{align}
I_{n+1} - I_n &= \int_0^\pi \frac{\sin^2((n+1)x) - \sin^2(nx)}{\sin x} \, \mathrm{d}x = \int_0^\pi \frac{\sin x \cdot \sin((2n+1)x)}{\sin x} \, \mathrm{d}x \nonumber \\
&= \int_0^\pi \sin((2n+1)x) \, \mathrm{d}x = \frac{1}{2n+1} \cos((2n+1)x) \Big|_\pi^0 = \frac{2}{2n+1}.
\end{align}
于是
\[
\int_0^\pi \frac{\sin^2(nx)}{\sin x} \, \mathrm{d}x = I_n = \frac{2}{2n-1} + I_{n-1} = \cdots = \sum_{k=1}^{n-1} \frac{2}{2k+1} + I_1 = \sum_{k=0}^{n-1} \frac{2}{2k+1}=\sum_{k=1}^n{\frac{2}{2k-1}}.
\]
\end{enumerate}
\end{proof}

\begin{example}
设 \( a > 1 \)，计算积分 \( \int_{0}^{\frac{\pi}{2}} \ln(a^2 - \cos^2 x) \, \mathrm{d}x \)。
\end{example}
\begin{remark}
很多情况下不需求出被积函数的原函数，只需充分利用换元、分部积分以及被积函数的性质，即可求出积分的值。见下述解法二.
\end{remark}
\begin{solution}

{\color{blue}解法一:}设 \( a_0 = a > 1 \)。构造数列如下：
\[
a_{n + 1} = 2a_n^2 - 1 \quad (n = 0, 1, \cdots),
\]
则由\refexa{example:例3.82-利用双曲三角函数求递推数列}可知,存在 \( x_0 > 0 \) 使得
\[
a_0 = \text{ch}(x_0), \quad a_n = \text{ch}(2^n x_0),
\]
其中 \( \text{ch}(x) = \frac{1}{2}(\text{e}^x + \text{e}^{-x}) \)。可以解得
\begin{align}
x_0 = \ln \left( a_0 + \sqrt{a_0^2 - 1} \right). \label{5.1.16}
\end{align}
故
\[
a_n = \frac{\text{e}^{2^n x_0} + \text{e}^{-2^n x_0}}{2}.
\]
设
\[
I_n = \int_{0}^{\pi} \ln(a_n - \cos x) \, \mathrm{d}x,
\]
则
\begin{align*}
I_0 &= \int_{0}^{\pi} \ln(a_0 - \cos x) \, \mathrm{d}x = \int_{0}^{\frac{\pi}{2}} \ln(a_0 - \cos x) \, \mathrm{d}x + \int_{\frac{\pi}{2}}^{\pi} \ln(a_0 - \cos x) \, \mathrm{d}x
\\
&= \int_{0}^{\frac{\pi}{2}} \ln(a_0 - \cos x) \, \mathrm{d}x + \int_{0}^{\frac{\pi}{2}} \ln(a_0 + \cos x) \, \mathrm{d}x
\\
&= \int_{0}^{\frac{\pi}{2}} \ln(a_0^2 - \cos^2 x) \, \mathrm{d}x = \int_{0}^{\frac{\pi}{2}} \ln \left( a_0^2 - \frac{1 + \cos 2x}{2} \right) \, \mathrm{d}x
\\
&= \int_{0}^{\frac{\pi}{2}} \ln \left( \frac{a_1 - \cos 2x}{2} \right) \, \mathrm{d}x = \frac{1}{2} \int_{0}^{\pi} \ln \left( \frac{a_1 - \cos x}{2} \right) \, \mathrm{d}x = \frac{1}{2} I_1 - \frac{\pi}{2} \ln 2.
\end{align*}
同理，有
\begin{align}
I_n = \frac{1}{2} I_{n + 1} - \frac{\pi}{2} \ln 2. \label{5.1.17}
\end{align}
由此递推公式，可得
\begin{align}
I_0 = \frac{1}{2^n} I_n - \left( 1 + \frac{1}{2} + \cdots + \frac{1}{2^{n - 1}} \right) \frac{\pi}{2} \ln 2. \label{5.1.18}
\end{align}
因为
\begin{align*}
I_n &= \int_{0}^{\pi} \ln(a_n - \cos x) \, \mathrm{d}x = \int_{0}^{\pi} \ln \left( \frac{\text{e}^{2^n x_0} + \text{e}^{-2^n x_0}}{2} - \cos x \right) \, \mathrm{d}x
\\
&= 2^n x_0 \pi + \int_{0}^{\pi} \ln \left( \frac{1 + \text{e}^{-2^{n + 1} x_0}}{2} - \text{e}^{-2^n x_0} \cos x \right) \, \mathrm{d}x,
\end{align*}
所以
\[
\frac{1}{2^n} I_n \to x_0 \pi \quad (n \to +\infty).
\]
故从式 \(\eqref{5.1.18}\) 可得
\[
I_0 = x_0 \pi - \pi \ln 2 = \pi \ln \left( \frac{a_0 + \sqrt{a_0^2 - 1}}{2} \right),
\]
即所求的积分为
\[
\int_{0}^{\frac{\pi}{2}} \ln(a^2 - \sin^2 x) \, \mathrm{d}x = \int_{0}^{\frac{\pi}{2}} \ln(a^2 - \cos^2 x) \, \mathrm{d}x = \pi \ln \left( \frac{a + \sqrt{a^2 - 1}}{2} \right).
\]

{\color{blue}解法二:}我们有
\[
F(a) = \int_{0}^{\frac{\pi}{2}} \ln(a^2 - \cos^2 x) \, \mathrm{d}x = \int_{0}^{\pi} \ln(a - \cos x) \, \mathrm{d}x.
\]
由\refthe{theorem:含参量积分的可微性1},关于 \( a \) 求导得到
\begin{align*}
F'(a) = \int_{0}^{\pi} \frac{1}{a - \cos x} \, \mathrm{d}x
\xlongequal{\text{万能公式}} \int_{0}^{+\infty} \frac{2}{a(1 + t^2) - (1 - t^2)} \, \mathrm{d}t = \frac{\pi}{\sqrt{a^2 - 1}}, \quad a > 1.
\end{align*}
因此
\[
F(a) =\int_1^a{F' \left( a \right) \mathrm{d}a}= \pi \ln \left( a + \sqrt{a^2 - 1} \right) + C, \quad a > 1.
\]
结合
\[
F(1)=2\int_0^{\frac{\pi}{2}}{\ln\sin x\,\mathrm{d}x}=-\pi \ln 2.
\]
可得
\[
\int_{0}^{\frac{\pi}{2}} \ln(a^2 - \cos^2 x) \, \mathrm{d}x = \pi \ln \left( \frac{a + \sqrt{a^2 - 1}}{2} \right), \quad a > 1.
\]

\end{solution}

\begin{proposition}[周期函数在任意一个周期上的定积分相同]\label{proposition:周期函数在任意一个周期上的定积分相同}
设$f$是$\mathbb{R}$上的可积函数且周期为$T>0$，则对$\forall a\in \mathbb{R}$，都有
\begin{align*}
\int_a^{a+T} f(x) \mathrm{d}x = \int_0^T f(x) \mathrm{d}x.
\end{align*}
\end{proposition}
\begin{proof}
由条件可得
\begin{align*}
\int_a^{a+T} f(x) \mathrm{d}x &= \int_a^{\left( \lfloor \frac{a}{T} \rfloor +1 \right) T} f(x) \mathrm{d}x + \int_{\left( \lfloor \frac{a}{T} \rfloor +1 \right) T}^{a+T} f(x) \mathrm{d}x 
= \int_a^{\left( \lfloor \frac{a}{T} \rfloor +1 \right) T} f(x) \mathrm{d}x + \int_{\lfloor \frac{a}{T} \rfloor T}^a f(x+T) \mathrm{d}x \\
&= \int_a^{\left( \lfloor \frac{a}{T} \rfloor +1 \right) T} f(x) \mathrm{d}x + \int_{\lfloor \frac{a}{T} \rfloor T}^a f(x) \mathrm{d}x 
= \int_{\lfloor \frac{a}{T} \rfloor T}^{\left( \lfloor \frac{a}{T} \rfloor +1 \right) T} f(x) \mathrm{d}x \\
&= \int_0^T f\left( x + \lfloor \frac{a}{T} \rfloor T \right) \mathrm{d}x= \int_0^T f(x) \mathrm{d}x.
\end{align*}
\end{proof}

\begin{theorem}[点火公式]\label{theorem:点火公式}
\begin{enumerate}
\item 记$I_n=\int_0^{\frac{\pi}{2}}\sin^n x\mathrm{d}x=\int_0^{\frac{\pi}{2}}\cos^n x\mathrm{d}x,\forall n\in \mathbb{N}$,则
\begin{align*}
I_n=\frac{n-1}{n}I_{n-2},\quad \forall n\geqslant 2.
\end{align*}
从而
\begin{align}\label{eq:点火公式1}
I_n=
\begin{cases}
\frac{(n-1)!!}{n!!}I_0=\frac{(n-1)!!}{n!!}\cdot \frac{\pi}{2}&,n\text{为偶数}\\
\frac{(n-1)!!}{n!!}I_1=\frac{(n-1)!!}{n!!}&,n\text{为奇数}
\end{cases}.
\end{align}

\item 记$J(m,n)=\int_0^{\frac{\pi}{2}}\sin^m x\cos^n x\mathrm{d}x,\forall n,m\in \mathbb{N}$,则
\begin{align*}
J(m,n)=\frac{m-1}{m+n}J(m-2,n),\quad \forall n,m\geqslant 2.
\end{align*}
\begin{align*}
J(m,n)=\frac{n-1}{m+n}J(m,n-2),\quad \forall n,m\geqslant 2.
\end{align*}
从而
\begin{align}\label{eq:点火公式2}
J(m,n)=
\begin{cases}
\frac{(m-1)!!(n-1)!!}{(m+n)!!}&,m,n\text{不全为偶数}\\
\frac{(m-1)!!(n-1)!!}{(m+n)!!}\cdot \frac{\pi}{2}&,m,n\text{全为偶数}
\end{cases}.
\end{align}

\end{enumerate}
\end{theorem}





\subsection{区间再现}

\begin{theorem}[区间再现恒等式]\label{theorem:区间再现恒等式}
当下述积分有意义时，我们有
\begin{enumerate}
\item \begin{align*}
\int_{a}^{b} f(x)\mathrm{d}x = \int_{a}^{b} f(a + b - x)\mathrm{d}x
=\frac{1}{2}\int_{a}^{b} [f(x)+f(a + b - x)]\mathrm{d}x
=\int_{a}^{\frac{a + b}{2}} [f(x)+f(a + b - x)]\mathrm{d}x.
\end{align*}

\item \begin{align*}
\int_{0}^{\infty} f(x)\mathrm{d}x = \int_{0}^{1} f(x)\mathrm{d}x+\int_{1}^{\infty} f(x)\mathrm{d}x
=\int_{0}^{1} \left[f(x)+\frac{f(\frac{1}{x})}{x^2}\right]\mathrm{d}x.
\end{align*}
\end{enumerate}
\end{theorem}
\begin{note}
注意:倒代换具有将$[0,1]$转化为$[1,+\infty)$的功能.
\end{note}
\begin{proof}
证明是显然的.(第1问中最后一个等号是由轴对称得到的)
\end{proof}

\begin{proposition}\label{example:常见积分1}
证明
\begin{enumerate}
\item \(\int_{0}^{\frac{\pi}{2}} \ln\sin x\mathrm{d}x=-\frac{\pi}{2}\ln 2.\)

\item \(\int_{0}^{\frac{\pi}{2}} \ln\cos x\mathrm{d}x=-\frac{\pi}{2}\ln 2.\)

\item \(\int_{0}^{1} \frac{\ln(1 + x)}{1 + x^2}\mathrm{d}x=\frac{\pi}{8}\ln 2.\)
\end{enumerate}
\end{proposition}
\begin{proof}
\begin{enumerate}
\item \begin{align*}
I&=\int_0^{\frac{\pi}{2}}{\ln\sin x\mathrm{d}x}=\int_0^{\frac{\pi}{4}}{\left[ \ln\sin x+\ln \sin \left( \frac{\pi}{2}-x \right) \right] \mathrm{d}x}
\\
&=\int_0^{\frac{\pi}{4}}{\ln \left( \cos x\sin x \right) \mathrm{d}x}=\int_0^{\frac{\pi}{4}}{\ln \frac{1}{2}\mathrm{d}x}+\int_0^{\frac{\pi}{4}}{\ln \left( \sin 2x \right) \mathrm{d}x}
\\
&=-\frac{\pi}{4}\ln 2+\frac{1}{2}\int_0^{\frac{\pi}{2}}{\ln\sin x\mathrm{d}x}=-\frac{\pi}{4}\ln 2+\frac{1}{2}I
\\
&\Longrightarrow I=\int_0^{\frac{\pi}{2}}{\ln\cos x\mathrm{d}x}=-\frac{\pi}{2}\ln 2.
\end{align*}

\item \begin{align*}
I&=\int_0^{\frac{\pi}{2}}{\ln\cos x\mathrm{d}x}=\int_0^{\frac{\pi}{2}}{\ln\sin x\mathrm{d}x}=\int_0^{\frac{\pi}{4}}{\left[ \ln\sin x+\ln \sin \left( \frac{\pi}{2}-x \right) \right] \mathrm{d}x}
\\
&=\int_0^{\frac{\pi}{4}}{\ln \left( \cos x\sin x \right) \mathrm{d}x}=\int_0^{\frac{\pi}{4}}{\ln \frac{1}{2}\mathrm{d}x}+\int_0^{\frac{\pi}{4}}{\ln \left( \sin 2x \right) \mathrm{d}x}
\\
&=-\frac{\pi}{4}\ln 2+\frac{1}{2}\int_0^{\frac{\pi}{2}}{\ln\sin x\mathrm{d}x}=-\frac{\pi}{4}\ln 2+\frac{1}{2}I
\\
&\Longrightarrow I=\int_0^{\frac{\pi}{2}}{\ln\cos x\mathrm{d}x}=-\frac{\pi}{2}\ln 2.
\end{align*}

\item {\color{blue}解法一:}
\begin{align*}
\int_0^1{\frac{\ln \left( 1+x \right)}{1+x^2}\mathrm{d}x}&\xlongequal{x=\tan \theta}\int_0^{\frac{\pi}{4}}{\frac{\ln \left( 1+\tan \theta \right)}{1+\tan \theta ^2}\mathrm{d}\tan \theta}=\int_0^{\frac{\pi}{4}}{\frac{\sec ^2\theta \cdot \ln \left( 1+\tan \theta \right)}{\sec ^2\theta}\mathrm{d}\theta}
\\
&=\int_0^{\frac{\pi}{4}}{\ln \left( 1+\tan \theta \right) \mathrm{d}\theta}=\int_0^{\frac{\pi}{8}}{\left[ \ln \left( 1+\tan \theta \right) +\ln \left( 1+\tan \left( \frac{\pi}{4}-\theta \right) \right) \right] \mathrm{d}\theta}
\\
&=\int_0^{\frac{\pi}{8}}{\left[ \ln \left( 1+\tan \theta \right) +\ln \left( 1+\frac{1-\tan \theta}{1+\tan \theta} \right) \right] \mathrm{d}\theta}
\\
&=\int_0^{\frac{\pi}{8}}{\left[ \ln \left( 1+\tan \theta \right) +\ln \frac{2}{1+\tan \theta} \right] \mathrm{d}\theta}
\\
&=\int_0^{\frac{\pi}{8}}{\ln 2\mathrm{d}\theta}=\frac{\pi}{8}\ln 2.
\end{align*} 

{\color{blue}解法二:}考虑含参量积分
\[
\varphi(\alpha) = \int_{0}^{1} \frac{\ln(1 + \alpha x)}{1 + x^2} \, \mathrm{d}x, \quad \alpha \in [0,1].
\]
显然 \( \varphi(0) = 0 \)，\( \varphi(1) = I \)，且函数 \( \frac{\ln(1 + \alpha x)}{1 + x^2} \) 在 \( R = [0,1] \times [0,1] \) 上满足\refthe{theorem:含参量积分的可微性1} 的条件，于是
\[
\varphi'(\alpha) = \int_{0}^{1} \frac{x}{(1 + x^2)(1 + \alpha x)} \, \mathrm{d}x.
\]
因为
\[
\frac{x}{(1 + x^2)(1 + \alpha x)} = \frac{1}{1 + \alpha^2} \left( \frac{\alpha + x}{1 + x^2} - \frac{\alpha}{1 + \alpha x} \right),
\]
所以
\begin{align*}
\varphi' \left( \alpha \right) &=\frac{1}{1+\alpha ^2}\left( \int_0^1{\frac{\alpha}{1+x^2}\,\mathrm{d}x}+\int_0^1{\frac{x}{1+x^2}\,\mathrm{d}x}-\int_0^1{\frac{\alpha}{1+\alpha x}\,\mathrm{d}x} \right) 
\\
&=\frac{1}{1+\alpha ^2}\left[ \alpha \arctan  x\Big|_{0}^{1}+\frac{1}{2}\ln \left( 1+x^2 \right) \Big|_{0}^{1}-\ln \left( 1+\alpha x \right) \Big|_{0}^{1} \right] 
\\
&=\frac{1}{1+\alpha ^2}\left[ \alpha \cdot \frac{\pi}{4}+\frac{1}{2}\ln 2-\ln \left( 1+\alpha \right) \right] .
\end{align*}
因此
\begin{align*}
\int_0^1{\varphi' \left( \alpha \right) \,\mathrm{d}\alpha}&=\int_0^1{\frac{1}{1+\alpha ^2}\left[ \frac{\pi}{4}\alpha +\frac{1}{2}\ln 2-\ln \left( 1+\alpha \right) \right] \,\mathrm{d}\alpha}
\\
&=\frac{\pi}{8}\ln \left( 1+\alpha ^2 \right) \Big|_{0}^{1}+\frac{1}{2}\ln 2\arctan  \alpha \Big|_{0}^{1}-\varphi \left( 1 \right) 
\\
&=\frac{\pi}{8}\ln 2+\frac{\pi}{8}\ln 2-\varphi \left( 1 \right) 
\\
&=\frac{\pi}{4}\ln 2-\varphi \left( 1 \right) .
\end{align*}
另一方面，
\[
\int_{0}^{1} \varphi'(\alpha) \, d\alpha = \varphi(1) - \varphi(0) = \varphi(1),
\]
所以 \( I = \varphi(1) = \frac{\pi}{8} \ln 2 \).
\end{enumerate}
\end{proof}

\begin{example}
计算
\begin{enumerate}
\item \(\int_{0}^{\infty} \frac{\ln x}{x^{2}+a^{2}}\mathrm{d}x\),$a>0$.

\item \(\int_{0}^{\infty} \frac{\ln x}{x^{2}+x + 1}\mathrm{d}x\).

\item\(\int_{0}^{1} \frac{\ln x}{\sqrt{x - x^{2}}}\mathrm{d}x\).
\end{enumerate}
\end{example}
\begin{solution}
\begin{enumerate}
\item 注意到
\begin{align}
\int_0^{+\infty}\frac{\ln x}{x^2+a^2}\mathrm{d}x\xlongequal{x=at}\frac{1}{a}\int_0^{+\infty}\frac{\ln (at)}{1+t^2}\mathrm{d}t
=\frac{1}{a}\int_0^{+\infty}\frac{\ln a}{1+t^2}\mathrm{d}t+\frac{1}{a}\int_0^{+\infty}\frac{\ln t}{1+t^2}\mathrm{d}t
=\frac{\pi \ln a}{2a}+\frac{1}{a}\int_0^{+\infty}\frac{\ln t}{1+t^2}\mathrm{d}t.\label{eq:1.1}
\end{align}
又注意到
\begin{align*}
\int_0^{+\infty}\frac{\ln t}{1+t^2}\mathrm{d}t\xlongequal{t=\frac{1}{x}}\int_0^{+\infty}\frac{\ln \frac{1}{x}}{1+\frac{1}{x^2}}\frac{1}{x^2}\mathrm{d}x=\int_0^{+\infty}\frac{-\ln x}{1+x^2}\mathrm{d}x\Longrightarrow \int_0^{+\infty}\frac{\ln t}{1+t^2}\mathrm{d}t=0.
\end{align*}
于是再结合\eqref{eq:1.1}式可得
\begin{align*}
\int_0^{+\infty}\frac{\ln x}{x^2+a^2}\mathrm{d}x=\frac{\pi \ln a}{2a}.
\end{align*}

\item \begin{align*}
\int_0^{\infty}{\frac{\ln x}{x^2+x+1}\mathrm{d}x}\xlongequal{x=\frac{1}{t}}\int_0^{\infty}{\frac{-\ln t}{1+\frac{1}{t}+\frac{1}{t^2}}\mathrm{d}\frac{1}{t}}=\int_0^{+\infty}{\frac{-\ln t}{1+t+t^2}\mathrm{d}t}\Longrightarrow \int_0^{\infty}{\frac{\ln x}{x^2+x+1}\mathrm{d}x}=0.
\end{align*}

\item\begin{align*}
\int_{0}^{1} \frac{\ln x}{\sqrt{x - x^{2}}}\mathrm{d}x &\xlongequal{x = \sin^{2}y} \int_{0}^{\frac{\pi}{2}} \frac{\ln \sin^{2}y}{\sqrt{\sin^{2}y(1 - \sin^{2}y)}}\mathrm{d}\sin^{2}y\\
&= 4\int_{0}^{\frac{\pi}{2}} \ln \sin y\mathrm{d}y \xlongequal{\text{\refpro{example:常见积分1}}}4\cdot\left(-\frac{\pi}{2}\ln 2\right)= -2\pi\ln 2.
\end{align*} 
\end{enumerate}

\end{solution}

\begin{example}
\begin{enumerate}
\item 对\(n \in \mathbb{N}\)，计算\(\int_{-\pi}^{\pi} \frac{\sin(nx)}{(1 + 2^{x})\sin x}\mathrm{d}x\)。

\item \(\int_{-\pi}^{\pi} \frac{x\sin x\arctan e^{x}}{1 + \cos^{2}x}\mathrm{d}x\)。

\item 对\(n \in \mathbb{N}\)，计算\(\int_{0}^{2\pi} \sin(\sin x + nx)\mathrm{d}x\).
\end{enumerate}
\end{example}
\begin{solution}
\begin{enumerate}
\item \begin{align*}
&\int_{-\pi}^{\pi}{\frac{\sin \left( nx \right)}{\left( 1+2^x \right) \sin x}\mathrm{d}x}=\int_{-\pi}^0{\left[ \frac{\sin \left( nx \right)}{\left( 1+2^x \right) \sin x}+\frac{\sin \left( nx \right)}{\left( 1+2^{-x} \right) \sin x} \right] \mathrm{d}x}=\int_{-\pi}^0{\frac{\sin \left( nx \right)}{\sin x}\left( \frac{1}{1+2^x}+\frac{1}{1+2^{-x}} \right) \mathrm{d}x}
\\
&=\int_{-\pi}^0{\frac{\sin \left( nx \right)}{\sin x}\cdot \frac{2+2^x+2^{-x}}{2+2^x+2^{-x}}\mathrm{d}x}=\int_{-\pi}^0{\frac{\sin \left( nx \right)}{\sin x}\mathrm{d}x}
=\int_0^{\pi}{\frac{\sin \left( nx \right)}{\sin x}\mathrm{d}x}\xlongequal{\text{\refexa{example:经典定积分(必记)}}}\begin{cases}
0,n\text{为偶数}\\
\pi ,n\text{为奇数}\\
\end{cases}.
\end{align*}

\item \begin{align*}
\int_{-\pi}^{\pi}{\frac{x\sin x\arctan  e^x}{1+\cos ^2x}\mathrm{d}x}&=\int_{-\pi}^0{\left( \frac{x\sin x\arctan  e^x}{1+\cos ^2x}+\frac{x\sin x\arctan  e^{-x}}{1+\cos ^2x} \right) \mathrm{d}x}=\int_{-\pi}^0{\frac{x\sin x}{1+\cos ^2x}\left( \arctan  e^x+\arctan  e^{-x} \right) \mathrm{d}x}
\\
&\xlongequal{\text{\nrefpro{proposition:arctan相关等式}{(1)}}}\int_{-\pi}^0{\frac{x\sin x}{1+\cos ^2x}\cdot \frac{\pi}{2}\mathrm{d}x}=\frac{\pi}{2}\int_0^{\pi}{\frac{x\sin x}{1+\cos ^2x}\mathrm{d}x}
\\
&=\frac{\pi}{2}\int_0^{\frac{\pi}{2}}{\left( \frac{x\sin x}{1+\cos ^2x}+\frac{\left( \pi -x \right) \sin x}{1+\cos ^2x} \right) \mathrm{d}x}=\frac{\pi ^2}{2}\int_0^{\frac{\pi}{2}}{\frac{\sin x}{1+\cos ^2x}\mathrm{d}x}
\\
&=\frac{\pi ^2}{2}\mathrm{arctan}\cos x\Big |_{\frac{\pi}{2}}^{0}=\frac{\pi ^2}{2}\cdot \frac{\pi}{4}=\frac{\pi ^3}{8}.
\end{align*}

\item\begin{align*}
\int_0^{2\pi}{\sin \left( \sin x+nx \right) \mathrm{d}x}&=\int_0^{2\pi}{\sin \left[ \sin \left( 2\pi -x \right) +n\left( 2\pi -x \right) \right] \mathrm{d}x}
\\
&=\int_0^{2\pi}{\sin \left( -\sin x-nx \right) \mathrm{d}x}=-\int_0^{2\pi}{\sin \left( \sin x+nx \right) \mathrm{d}x}
\\
&\Longrightarrow \int_0^{2\pi}{\sin \left( \sin x+nx \right) \mathrm{d}x}=0.
\end{align*} 
\end{enumerate}

\end{solution}

\begin{example}
计算$\int_{0}^{+\infty} \frac{\mathrm{d}x}{\left(1 + x^2\right)\left(1 + x^{2019}\right)}$。
\end{example}
\begin{solution}
\begin{align*}
&\int_0^{+\infty}{\frac{\mathrm{d}x}{\left( 1+x^2 \right) \left( 1+x^{2019} \right)}}=\int_0^{+\infty}{\frac{\mathrm{d}x}{\left( 1+x^2 \right) \left( 1+x^{2019} \right)}}
\\
&\xlongequal{t=\frac{1}{x}}\int_0^{+\infty}{\frac{t^{2019}\mathrm{d}t}{\left( 1+t^2 \right) \left( 1+t^{2019} \right)}}=\frac{1}{2}\int_0^{+\infty}{\frac{\left( 1+x^{2019} \right) \mathrm{d}x}{\left( 1+x^2 \right) \left( 1+x^{2019} \right)}}
\\
&=\frac{1}{2}\int_0^{+\infty}{\frac{\mathrm{d}x}{1+x^2}}=\frac{\pi}{4}.
\end{align*}

\end{solution}




\subsection{Frullani(傅汝兰尼)积分}

\begin{theorem}[Frullani(傅汝兰尼)积分]\label{theorem:Frullani(傅汝兰尼)积分}
设 $f \in C(0,+\infty)$.
\begin{enumerate}
\item 若存在极限
\begin{align}
\lim_{x \to 0^+} f(x), \lim_{x \to +\infty} f(x), \label{equation:::::::-----2864611---14.3}
\end{align}
则对 $a,b > 0$ 有
$$\int_0^\infty \frac{f(ax) - f(bx)}{x} \mathrm{d}x = \left[ \lim_{x \to 0^+} f(x) - \lim_{x \to +\infty} f(x) \right] \ln \frac{b}{a}.$$

\item 若存在极限和积分
\begin{align}
\lim_{x \to 0^+} f(x) = \alpha, \int_A^\infty \frac{f(x)}{x} \mathrm{d}x. \label{equation:::::::-----2864611---14.4}
\end{align}
则对 $a,b > 0$, 有
$$\int_0^\infty \frac{f(ax) - f(bx)}{x} \mathrm{d}x = \alpha \ln \frac{b}{a}.$$

\item 若存在极限和积分
\begin{align}
\lim_{x \to +\infty} f(x) = \alpha, \int_0^1 \frac{f(x)}{x} \mathrm{d}x. \label{equation:::::::-----2864611---14.5}
\end{align}
则对 $a,b > 0$, 有
$$\int_0^\infty \frac{f(ax) - f(bx)}{x} \mathrm{d}x = \alpha \ln \frac{a}{b}.$$

\item 若 $f$ 是周期 $T > 0$ 函数且 $\lim_{x \to 0^+} f(x)$ 存在, 则对 $a,b > 0$ 有
$$\int_0^\infty \frac{f(ax) - f(bx)}{x} \mathrm{d}x = \left[ \lim_{x \to 0^+} f(x) - \frac{1}{T} \int_0^T f(x) \mathrm{d}x \right] \ln \frac{b}{a}.$$

\item 若 $f$ 满足 $\lim_{x \to 0^+} f(x), \lim_{x \to +\infty} \frac{1}{x} \int_0^x f(y) \mathrm{d}y$ 存在, 则对 $a,b > 0$ 有
$$\int_0^\infty \frac{f(ax) - f(bx)}{x} \mathrm{d}x = \left[ \lim_{x \to 0^+} f(x) - \lim_{x \to +\infty} \frac{1}{x} \int_0^x f(y) \mathrm{d}y \right] \ln \frac{b}{a}.$$
\end{enumerate}
\end{theorem}
\begin{note}
傅汝兰尼积分有诸多变种, 无需记忆具体表达式, 知道有大概这么一个东西即可.
\end{note}
\begin{proof}
不妨设 $b > a$.
\begin{enumerate}
\item 给定 $A > \delta > 0$, 考虑
\begin{align*}
\int_\delta^A \frac{f(ax) - f(bx)}{x} \mathrm{d}x &= \int_\delta^A \frac{f(ax)}{x} \mathrm{d}x - \int_\delta^A \frac{f(bx)}{x} \mathrm{d}x \\
&= \int_{a\delta}^{aA} \frac{f(x)}{x} \mathrm{d}x - \int_{b\delta}^{bA} \frac{f(x)}{x} \mathrm{d}x \\
&= \int_{bA}^{aA} \frac{f(x)}{x} \mathrm{d}x - \int_{b\delta}^{a\delta} \frac{f(x)}{x} \mathrm{d}x \\
&\xlongequal{\text{\hyperref[theorem:积分中值定理]{积分中值定理}}} f(\theta_1) \int_{bA}^{aA} \frac{1}{x} \mathrm{d}x - f(\theta_2) \int_{b\delta}^{a\delta} \frac{1}{x} \mathrm{d}x,
\end{align*}
这里 $\theta_1 \in (aA, bA), \theta_2 \in (a\delta, b\delta)$, 于是让 $A \to +\infty, \delta \to 0^+$, 由\eqref{equation:::::::-----2864611---14.3}, 我们知
$$\int_0^\infty \frac{f(ax) - f(bx)}{x} \mathrm{d}x = \left[ \lim_{x \to 0^+} f(x) - \lim_{x \to +\infty} f(x) \right] \ln \frac{b}{a}.$$

\item 给定 $A > \delta > 0$, 考虑
\begin{align*}
\int_\delta^A \frac{f(ax) - f(bx)}{x} \mathrm{d}x &= \int_\delta^A \frac{f(ax)}{x} \mathrm{d}x - \int_\delta^A \frac{f(bx)}{x} \mathrm{d}x \\
&= \int_{a\delta}^{aA} \frac{f(x)}{x} \mathrm{d}x - \int_{b\delta}^{bA} \frac{f(x)}{x} \mathrm{d}x \\
&= \int_{bA}^{aA} \frac{f(x)}{x} \mathrm{d}x - \int_{b\delta}^{a\delta} \frac{f(x)}{x} \mathrm{d}x \\
&\xlongequal{\text{\hyperref[theorem:积分中值定理]{积分中值定理}}} \int_{bA}^{aA} \frac{f(x)}{x} \mathrm{d}x - f(\theta) \int_{b\delta}^{a\delta} \frac{1}{x} \mathrm{d}x,
\end{align*}
这里 $\theta \in (a\delta, b\delta)$, 于是让 $A \to +\infty, \delta \to 0^+$, 由 \eqref{equation:::::::-----2864611---14.4}, 我们知
$$\int_0^\infty \frac{f(ax) - f(bx)}{x} \mathrm{d}x = \alpha \ln \frac{b}{a}.$$

\item 给定 $A > \delta > 0$, 考虑
\begin{align*}
\int_\delta^A \frac{f(ax) - f(bx)}{x} \mathrm{d}x &= \int_\delta^A \frac{f(ax)}{x} \mathrm{d}x - \int_\delta^A \frac{f(bx)}{x} \mathrm{d}x \\
&= \int_{a\delta}^{aA} \frac{f(x)}{x} \mathrm{d}x - \int_{b\delta}^{bA} \frac{f(x)}{x} \mathrm{d}x \\
&= \int_{bA}^{aA} \frac{f(x)}{x} \mathrm{d}x - \int_{b\delta}^{a\delta} \frac{f(x)}{x} \mathrm{d}x \\
&\xlongequal{\text{\hyperref[theorem:积分中值定理]{积分中值定理}}} f(\theta) \int_{bA}^{aA} \frac{1}{x} \mathrm{d}x - \int_{b\delta}^{a\delta} \frac{f(x)}{x} \mathrm{d}x,
\end{align*}
这里 $\theta \in (aA, bA)$, 于是让 $A \to +\infty, \delta \to 0^+$, 由\eqref{equation:::::::-----2864611---14.5}, 我们知
$$\int_0^\infty \frac{f(ax) - f(bx)}{x} \mathrm{d}x = \alpha \ln \frac{a}{b}.$$

\item 给定 $A > \delta > 0$, 考虑
\begin{align*}
\int_\delta^A \frac{f(ax) - f(bx)}{x} \mathrm{d}x &= \int_\delta^A \frac{f(ax)}{x} \mathrm{d}x - \int_\delta^A \frac{f(bx)}{x} \mathrm{d}x \\
&= \int_{a\delta}^{aA} \frac{f(x)}{x} \mathrm{d}x - \int_{b\delta}^{bA} \frac{f(x)}{x} \mathrm{d}x \\
&= \int_{bA}^{aA} \frac{f(x)}{x} \mathrm{d}x - \int_{b\delta}^{a\delta} \frac{f(x)}{x} \mathrm{d}x \\
&\xlongequal{\text{\hyperref[theorem:积分中值定理]{积分中值定理}}} \int_{bA}^{aA} \frac{f(x)}{x} \mathrm{d}x - f(\theta) \int_{b\delta}^{a\delta} \frac{1}{x} \mathrm{d}x \\
&= \int_b^a \frac{f(Ax)}{x} \mathrm{d}x - f(\theta) \int_{b\delta}^{a\delta} \frac{1}{x} \mathrm{d}x,
\end{align*}
这里 $\theta \in (a\delta, b\delta)$. 现在
$$\lim_{\delta \to 0^+} \left( -f(\theta) \int_{b\delta}^{a\delta} \frac{1}{x} \mathrm{d}x \right) = \lim_{x \to 0^+} f(x) \ln \frac{b}{a}.$$

由\hyperref[theorem:Riemann引理]{Riemann引理}, 我们有
$$\lim_{A \to +\infty} \int_b^a \frac{f(Ax)}{x} \mathrm{d}x = \int_b^a \frac{1}{x} \mathrm{d}x \cdot \frac{1}{T} \int_0^T f(x) \mathrm{d}x = -\frac{1}{T} \int_0^T f(x) \mathrm{d}x \cdot \ln \frac{b}{a},$$
这就证明了
$$\int_0^\infty \frac{f(ax) - f(bx)}{x} \mathrm{d}x = \left[ \lim_{x \to 0^+} f(x) - \frac{1}{T} \int_0^T f(x) \mathrm{d}x \right] \ln \frac{b}{a}.$$

\item 上一问证明中把使用的\hyperref[theorem:Riemann引理]{Riemann引理}用\hyperref[theorem:L^1版本的Riemann引理]{平均值极限版本的Riemann引理}代替即可.
\end{enumerate}
\end{proof}





\subsection{化成多元累次积分(换序)}

\begin{proposition}[重要定积分结论]\label{proposition:重要定积分结论}
证明:
\begin{enumerate}[(1)]
\item\label{proposition:重要定积分结论-1} $\mathbf{Gauss-Poisson}$\textbf{(高斯-泊松)积分}:\(\int_{0}^{\infty} e^{-x^{2}}\mathrm{d}x=\frac{\sqrt{\pi}}{2}\).

\item\label{proposition:重要定积分结论-2} \(\int_{0}^{+\infty} \frac{\sin x}{x}\mathrm{d}x=\frac{\pi}{2}\),进而$\int_{-\infty}^{+\infty} \frac{\sin x}{x}\mathrm{d}x=\pi$.

\item\label{proposition:重要定积分结论-3} $\int_0^{+\infty}{\frac{\sin ^2x}{x^2}\mathrm{d}x}=\frac{\pi}{2}.$

\item\label{proposition:重要定积分结论-4} $\mathbf{Fresnel}$\textbf{(菲涅尔)积分}:\(\int_{0}^{\infty} \sin x^{2}\mathrm{d}x=\int_{0}^{\infty} \cos x^{2}\mathrm{d}x=\sqrt{\frac{\pi}{8}}\). 
\end{enumerate}
\end{proposition}
\begin{proof}
\begin{enumerate}[(1)]
\item 注意到
\begin{align*}
\left( \int_0^{+\infty}{e^{-x^2}\mathrm{d}x} \right) ^2&=\left( \int_0^{+\infty}{e^{-x^2}\mathrm{d}x} \right) \left( \int_0^{+\infty}{e^{-y^2}\mathrm{d}y} \right)=\int_0^{+\infty}{e^{-x^2}\left( \int_0^{+\infty}{e^{-y^2}}\mathrm{d}x \right) \mathrm{d}y}
\\
&=\int_0^{+\infty}{\left( \int_0^{+\infty}{e^{-\left( x^2+y^2 \right)}}\mathrm{d}x \right) \mathrm{d}y}\xlongequal{e^{-\left( x^2+y^2 \right)}\text{连续}}\iint_{R^2}{e^{-\left( x^2+y^2 \right)}\mathrm{d}x\mathrm{d}y}
\\
&=\int_0^{\frac{\pi}{2}}{\mathrm{d}\theta}\int_0^{+\infty}{re^{-r^2}\mathrm{d}r}=\frac{\pi}{2}\int_0^{+\infty}{re^{-r^2}\mathrm{d}r}
\\
&=\frac{\pi}{4}\int_0^{+\infty}{e^{-r^2}\mathrm{d}r^2}=\frac{\pi}{4}.
\end{align*}
故\(\int_{0}^{\infty} e^{-x^{2}}\mathrm{d}x=\frac{\sqrt{\pi}}{2}\).

\item 由分部积分可得
\begin{align*}
&\int_0^{+\infty}{e^{-yx}\sin x\,\mathrm{d}x}=e^{-yx}\cos x\big |_{x=+\infty}^{x=0}-y\int_0^{+\infty}{e^{-yx}\cos x\,\mathrm{d}x}=1-y\int_0^{+\infty}{e^{-yx}\cos x\,\mathrm{d}x}
\\
&=1-ye^{-yx}\sin x\big |_{x=+\infty}^{x=0}-y^2\int_0^{+\infty}{e^{-yx}\sin x\,\mathrm{d}x}=1-y^2\int_0^{+\infty}{e^{-yx}\sin x\,\mathrm{d}x}.
\end{align*}
因此
\begin{align*}
\int_0^{+\infty}{e^{-yx}\sin x\,\mathrm{d}x}=\frac{1}{1+y^2}.
\end{align*}
于是就有
\begin{align*}
\int_0^{+\infty}{\frac{\sin x}{x}\,\mathrm{d}x}
= \int_0^{+\infty}{\sin x \left( \int_0^{+\infty}{e^{-yx}\,\mathrm{d}y} \right) \,\mathrm{d}x} 
= \int_0^{+\infty}{\mathrm{d}y} \int_0^{+\infty}{e^{-yx}\sin x \,\mathrm{d}x} 
= \int_0^{+\infty}{\frac{1}{y^2 + 1}\,\mathrm{d}y} 
= \frac{\pi}{2}.
\end{align*}

\item 由分部积分可得
\begin{align*}
\int_0^{+\infty}{\frac{\sin ^2x}{x^2}\mathrm{d}x}&=\frac{1}{2}\int_0^{+\infty}{\frac{1-\cos 2x}{x^2}\mathrm{d}x}=-\frac{1}{2}\int_0^{+\infty}{\left( 1-\cos 2x \right) \mathrm{d}\frac{1}{x}}
\\
&=\frac{1}{2}\frac{1-\cos x}{x^2}\Big|_{+\infty}^{0}+\frac{1}{2}\int_0^{+\infty}{\frac{2\sin 2x}{x}\mathrm{d}x}
\\
&=\int_0^{+\infty}{\frac{\sin x}{x}\mathrm{d}x}\xlongequal{\text{\nrefpro{proposition:重要定积分结论}{(2)}}}\frac{\pi}{2}.
\end{align*}

\item 见\hyperref[Complex Analysis-proposition:几个特殊积分的计算-Fresnel积分]{Fresnel积分}.
\end{enumerate}
\end{proof}

\begin{proposition}
设$n,m\in \mathbb{N},\,n\geqslant m,$且$n-m$为偶数,计算
$\int_0^{+\infty}{\frac{(\sin x)^n}{x^m}\mathrm{d}x}.$
\end{proposition}
\begin{solution}
\begin{enumerate}[(1)]
\item 当$m=n=1$时,由\refpro{proposition:重要定积分结论}知$\int_0^{+\infty}{\frac{\sin x}{x}\mathrm{d}x}=\frac{\pi}{2}.$

\item 当$m=1,n=2k-1$时,由\refthe{theorem:sinnx和cosnx以及sin^nx和cos^nx的展开式}和\refpro{proposition:组合数恒等式1111}知
\begin{align*}
\int_0^{+\infty}{\frac{(\sin x)^{2k-1}}{x}\mathrm{d}x}&\xlongequal{\text{\refthe{theorem:sinnx和cosnx以及sin^nx和cos^nx的展开式}}}\frac{1}{2^{2k-2}}\sum\limits_{i=0}^{k-1}{(-1)^{k+i-1}\binom{2k-1}{i}\int_0^{+\infty}{\frac{\sin\left((2k-2i-1)x\right)}{x}\mathrm{d}x}}\\
&\xlongequal{y=(2k-2i-1)x}\frac{1}{2^{2k-2}}\sum\limits_{i=0}^{k-1}{(-1)^{k+i-1}\binom{2k-1}{i}\int_0^{+\infty}{\frac{\sin y}{y}\mathrm{d}y}}\\
&=\frac{\pi}{2^{2k-1}}\sum\limits_{i=0}^{k-1}{(-1)^{k+i-1}\binom{2k-1}{i}}=(-1)^{k+1}\frac{\pi}{2^{2k-1}}\sum\limits_{i=0}^{k-1}{(-1)^i\binom{2k-1}{i}}\\
&\xlongequal{\text{\refpro{proposition:组合数恒等式1111}}}\frac{\pi}{2^{2k-1}}\binom{2k-2}{k-1}=\frac{\pi}{2}\frac{(2k-2)!}{2^{k-1}(k-1)!\cdot 2^{k-1}(k-1)!}\\
&=\frac{\pi}{2}\frac{(2k-2)!!(2k-3)!!}{(2k-2)!!\cdot (2k-2)!!}=\frac{(2k-3)!!}{(2k-2)!!}\cdot \frac{\pi}{2}.
\end{align*}

\item 当$m=2,n=2k$时,利用分部积分得
\begin{align*}
&\int_0^{+\infty}{\frac{(\sin x)^{2k}}{x^2}\mathrm{d}x}=-\frac{(\sin x)^{2k}}{x}\bigg|_{0}^{+\infty}+2k\int_0^{+\infty}{\frac{(\sin x)^{2k-1}\cos x}{x}\mathrm{d}x}=2k\int_0^{+\infty}{\frac{(\sin x)^{2k-1}\cos x}{x}\mathrm{d}x}\\
&\xlongequal{\text{\refthe{theorem:sinnx和cosnx以及sin^nx和cos^nx的展开式}}}\frac{k}{2^{2k-3}}\sum\limits_{i=0}^{k-1}{(-1)^{k+i-1}\binom{2k-1}{i}\int_0^{+\infty}{\frac{\sin\left((2k-2i-1)x\right)\cos x}{x}\mathrm{d}x}}\\
&=\frac{k}{2^{2k-2}}\sum\limits_{i=0}^{k-1}{(-1)^{k+i-1}\binom{2k-1}{i}\int_0^{+\infty}{\frac{\left[\sin\left((2k-2i)x\right)+\sin\left((2k-2i-2)x\right)\right]}{x}\mathrm{d}x}}\\
&=\frac{k}{2^{2k-2}}\sum\limits_{i=0}^{k-1}{(-1)^{k+i-1}\binom{2k-1}{i}\int_0^{+\infty}{\frac{\sin\left((2k-2i)x\right)}{x}\mathrm{d}x}}+\frac{k}{2^{2k-2}}\sum\limits_{i=0}^{k-1}{(-1)^{k+i-1}\binom{2k-1}{i}\int_0^{+\infty}{\frac{\sin\left((2k-2i-2)x\right)}{x}\mathrm{d}x}}\\
&=\frac{(-1)^{k+1}k}{2^{2k-2}}\sum\limits_{i=0}^{k-1}{(-1)^i\binom{2k-1}{i}\int_0^{+\infty}{\frac{\sin x}{x}\mathrm{d}x}}+\frac{(-1)^{k+1}k}{2^{2k-2}}\sum\limits_{i=0}^{k-1}{(-1)^i\binom{2k-1}{i}\int_0^{+\infty}{\frac{\sin x}{x}\mathrm{d}x}}\\
&=\frac{\pi}{2}\cdot \frac{(-1)^{k+1}k}{2^{2k-2}}\sum\limits_{i=0}^{k-1}{(-1)^i\binom{2k-1}{i}}\xlongequal{\text{\refpro{proposition:组合数恒等式1111}}}\frac{\pi}{2}\cdot \frac{k}{2^{2k-2}}\binom{2k-2}{k-1}\\
&=\frac{\pi}{2}\cdot \frac{(2k-2)!}{2^{k-1}(k-1)!\cdot 2^{k-1}(k-1)!}=\frac{\pi}{2}\cdot \frac{(2k-2)!!(2k-3)!!}{(2k-2)!!\cdot (2k-2)!!}\\
&=\frac{(2k-3)!!}{(2k-2)!!}\cdot \frac{\pi}{2}.
\end{align*}

\item 当$m>2$时,利用分部积分得
\begin{align*}
\int_0^{+\infty}{\frac{(\sin x)^n}{x^m}\mathrm{d}x}&=-\frac{1}{m-1}\left[\frac{(\sin x)^n}{x^{m-1}}\bigg|_{0}^{+\infty}-n\int_0^{+\infty}{\frac{(\sin x)^{n-1}\cos x}{x^{m-1}}\mathrm{d}x}\right]\\
&=\frac{n}{m-1}\int_0^{+\infty}{\frac{(\sin x)^{n-1}\cos x}{x^{m-1}}\mathrm{d}x}=\frac{n}{(m-1)(m-2)}\int_0^{+\infty}{\frac{\left[(\sin x)^{n-1}\cos x\right]'}{x^{m-1}}\mathrm{d}x}\\
&=\frac{n(n-1)}{(m-1)(m-2)}\int_0^{+\infty}{\frac{(\sin x)^{n-2}\cos^2x}{x^{m-2}}\mathrm{d}x}-\frac{n}{(m-1)(m-2)}\int_0^{+\infty}{\frac{(\sin x)^n}{x^{m-2}}\mathrm{d}x}\\
&=\frac{n(n-1)}{(m-1)(m-2)}\int_0^{+\infty}{\frac{(\sin x)^{n-2}\left(1-\sin^2x\right)}{x^{m-2}}\mathrm{d}x}-\frac{n}{(m-1)(m-2)}\int_0^{+\infty}{\frac{(\sin x)^n}{x^{m-2}}\mathrm{d}x}\\
&=\frac{n(n-1)}{(m-1)(m-2)}\int_0^{+\infty}{\frac{(\sin x)^{n-2}}{x^{m-2}}\mathrm{d}x}-\frac{n^2}{(m-1)(m-2)}\int_0^{+\infty}{\frac{(\sin x)^n}{x^{m-2}}\mathrm{d}x}.
\end{align*}
不断分部积分,最后只需计算$m=1,n=2k-1$和$m=2,n=2k\left(k\in \mathbb{Z}_+\right)$时的情况.再利用(2)(3)的结论即得.
\end{enumerate}

\end{solution}


\begin{example}
计算 \(\int_{0}^{1}\sin\ln\frac{1}{x}\cdot\frac{x^{b}-x^{a}}{\ln x}\mathrm{d}x\ (b > a > 0)\)。
\end{example}
\begin{proof}
\begin{align*}
\int_0^1{\mathrm{sin}\ln \frac{1}{x}}\cdot \frac{x^b-x^a}{\ln x}\mathrm{d}x&=\int_0^1{\mathrm{sin}\ln \frac{1}{x}}\left( \int_a^b{x^y}\mathrm{d}y \right) \mathrm{d}x=\int_a^b{\mathrm{d}y}\int_0^1{x^y\mathrm{sin}\ln \frac{1}{x}}\mathrm{d}x
\\
&\xlongequal{x=e^{-t}}\int_a^b{\mathrm{d}y}\int_{+\infty}^0{e^{-ty}\sin t}\mathrm{d}e^{-t}=\int_a^b{\mathrm{d}y}\int_0^{+\infty}{e^{-t\left( y+1 \right)}\sin t}\mathrm{d}t
\\
&\xlongequal{\text{\nrefpro{proposition:重要定积分结论}{(2)的证明过程}}}\int_a^b{\frac{1}{1+\left( y+1 \right) ^2}\mathrm{d}y}=\arctan  \left( b+1 \right) -\arctan  \left( a+1 \right) .
\end{align*} 
\end{proof}


\subsection{化成含参积分(费曼积分法)}

\begin{example}
设\(a,b\geqslant 0\)且不全为\(0\)，计算\(\int_{0}^{\frac{\pi}{2}} \ln\left( a^{2} \cos^{2}x + b^{2} \sin^{2}x\right) \mathrm{d}x\).
\end{example}
\begin{remark}
实际上,根据$a>b$时得到的结果,可以看出$F(a,b)=\pi \ln\frac{a+b}{2}$对$a,b$有轮换对称性,故这个结果对其他情况显然也成立.
\end{remark}
\begin{proof}
设\(F(a,b)=\int_{0}^{\frac{\pi}{2}} \ln\left( a^{2} \cos^{2}x + b^{2} \sin^{2}x\right) \mathrm{d}x\)，当\(a > b\)时，则
\\
\begin{align*}
\frac{\partial}{\partial b}F(a,b) &= \int_{0}^{\frac{\pi}{2}} \frac{\partial}{\partial b} \ln\left( a^{2} \cos^{2}x + b^{2} \sin^{2}x\right) \mathrm{d}x = \int_{0}^{\frac{\pi}{2}} \frac{2b\sin^{2}x}{a^{2} \cos^{2}x + b^{2} \sin^{2}x} \mathrm{d}x \\
&= \int_{0}^{\frac{\pi}{2}} \frac{2b\tan^{2}x}{a^{2} + b^{2}\tan^{2}x} \mathrm{d}x = \int_{0}^{+\infty} \frac{2bt^{2}}{(a^{2} + b^{2}t^{2})(1 + t^{2})} \mathrm{d}t \\
&= \frac{1}{a^{2}-b^{2}} \int_{0}^{+\infty} \left( \frac{2a^{2}b}{a^{2} + b^{2}t^{2}} - \frac{2b}{1 + t^{2}} \right) \mathrm{d}t \\
&= \frac{1}{a^{2}-b^{2}} \int_{0}^{+\infty} \frac{2a^{2}b}{a^{2} + b^{2}t^{2}} \mathrm{d}t - \frac{2b}{a^{2}-b^{2}} \int_{0}^{+\infty} \frac{1}{1 + t^{2}} \mathrm{d}t \\
&= \frac{2b}{a^{2}-b^{2}} \int_{0}^{+\infty} \frac{1}{1 + \left( \frac{b}{a}t \right)^{2}} \mathrm{d}t - \frac{b\pi}{a^{2}-b^{2}} \\
&= \frac{2b}{a^{2}-b^{2}} \cdot \frac{a}{b} \cdot \frac{\pi}{2} - \frac{b\pi}{a^{2}-b^{2}} = \frac{\pi}{a + b}.
\end{align*}
于是
\begin{align*}
F(a,b) &= F(a,0) + \int_{0}^{b} \frac{\partial}{\partial b^\prime} F(a,b^\prime) \mathrm{d}b^\prime = F(a,0) + \int_{0}^{b} \frac{\pi}{a + b^\prime} \mathrm{d}b^\prime \\
&= 2\int_{0}^{\frac{\pi}{2}} \ln(a\cos x) \mathrm{d}x + \pi \ln\frac{a + b}{a} \xlongequal{\text{\refexa{example:常见积分1}}} \pi \ln\frac{a + b}{2}.
\end{align*}
当\(a < b\)时，类似可得\(F(a,b)=\pi \ln\frac{a + b}{2}\)。当\(a = b\)时，\(F(a,b)=\int_{0}^{\frac{\pi}{2}} \ln a^{2} \mathrm{d}x = \pi \ln a = \pi \ln\frac{a + b}{2}\)。

综上，对\(\forall a,b\geqslant 0\)，都有\(F(a,b)=\pi \ln\frac{a + b}{2}\)。 
\end{proof}

\begin{example}
设$a>0,b>0$，计算积分$\int_0^{+\infty}\frac{\ln(a^2+x^2)}{b^2+x^2}\mathrm{d}x.$
\end{example}
\begin{proof}
记$f(t)=\int_0^{+\infty}\frac{\ln \left( a^2+tx^2 \right)}{b^2+x^2}\mathrm{d}x$，则
\begin{align*}
\frac{\partial}{\partial t}\frac{\ln \left( a^2+tx^2 \right)}{b^2+x^2}=\frac{x^2}{\left( b^2+x^2 \right) \left( a^2+tx^2 \right)}.
\end{align*}
对$\forall t\in \left[ 0,1 \right]$，显然$\frac{\ln \left( a^2+tx^2 \right)}{b^2+x^2},\frac{x^2}{\left( b^2+x^2 \right) \left( a^2+tx^2 \right)}\in C\left[ 0,1 \right]$，且
\begin{align*}
\int_0^1\frac{\ln \left( a^2+tx^2 \right)}{b^2+x^2}\mathrm{d}x\leqslant \int_0^1\frac{\ln \left( a^2+t \right)}{b^2}\mathrm{d}x<+\infty,
\end{align*}
\begin{align*}
\int_1^{+\infty}\frac{\ln \left( a^2+tx^2 \right)}{b^2+x^2}\mathrm{d}x=\int_1^{+\infty}\frac{\ln \left( \frac{a^2}{x^2}+t \right) +2\ln x}{b^2+x^2}\mathrm{d}x
\leqslant \ln \left( a^2+t \right) \int_1^{+\infty}\frac{1}{b^2+x^2}\mathrm{d}x+\int_1^{+\infty}\frac{2\ln x}{b^2+x^2}\mathrm{d}x<+\infty.
\end{align*}
故$f(t)=\int_0^{+\infty}\frac{\ln \left( a^2+tx^2 \right)}{b^2+x^2}\mathrm{d}x$在$\left( 0,1 \right)$上收敛. 注意到对$\forall \left[ c,d \right] \subset \left( 0,1 \right)$，有
\begin{align*}
\int_0^1\frac{\partial}{\partial t}\frac{\ln \left( a^2+tx^2 \right)}{b^2+x^2}\mathrm{d}x&=\int_0^{+\infty}\frac{x^2}{\left( b^2+x^2 \right) \left( a^2+tx^2 \right)}\mathrm{d}x\\
&=\frac{1}{a^2-tb^2}\int_0^{+\infty}\left( \frac{a^2}{a^2+tx^2}-\frac{b^2}{b^2+x^2} \right) \mathrm{d}x\\
&=\frac{1}{a^2-tb^2}\left[ \frac{a}{\sqrt{t}}\arctan \left( \frac{x\sqrt{t}}{a} \right) \bigg|_{0}^{+\infty}-b\arctan \frac{x}{b}\bigg|_{0}^{+\infty} \right] \\
&=\frac{\pi}{2\sqrt{t}\left( a+b\sqrt{t} \right)}\leqslant \frac{\pi}{2\sqrt{c}\left( a+b\sqrt{c} \right)}.
\end{align*}
故$\int_0^1\frac{\partial}{\partial t}\frac{\ln \left( a^2+tx^2 \right)}{b^2+x^2}\mathrm{d}x$在$\left( 0,1 \right)$上内闭一致收敛. 于是由\hyperref[theorem:含参反常积分的可微性]{含参量积分的可微性}知
\begin{align*}
f'(t)&=\frac{\mathrm{d}}{\mathrm{d}t}\int_0^{+\infty}\frac{\ln \left( a^2+tx^2 \right)}{b^2+x^2}\mathrm{d}x=\int_0^{+\infty}\frac{\partial}{\partial t}\frac{\ln \left( a^2+tx^2 \right)}{b^2+x^2}\mathrm{d}x\\
&=\int_0^{+\infty}\frac{x^2}{\left( b^2+x^2 \right) \left( a^2+tx^2 \right)}\mathrm{d}x
=\frac{1}{a^2-tb^2}\int_0^{+\infty}\left( \frac{a^2}{a^2+tx^2}-\frac{b^2}{b^2+x^2} \right) \mathrm{d}x\\
&=\frac{1}{a^2-tb^2}\left[ \frac{a}{\sqrt{t}}\arctan \left( \frac{x\sqrt{t}}{a} \right) \bigg|_{0}^{+\infty}-b\arctan \frac{x}{b}\bigg|_{0}^{+\infty} \right] \\
&=\frac{\pi}{2\sqrt{t}\left( a+b\sqrt{t} \right)}.
\end{align*}
故
\begin{align*}
\int_0^{+\infty}\frac{\ln \left( a^2+x^2 \right)}{b^2+x^2}\mathrm{d}x&=f(1)=\int_0^1f'(t)\mathrm{d}t+f(0)\\
&=\int_0^1\frac{\pi}{2\sqrt{t}\left( a+b\sqrt{t} \right)}\mathrm{d}t+\int_0^{+\infty}\frac{2\ln a}{b^2+x^2}\mathrm{d}x\\
&\xlongequal{x=\sqrt{t}}\pi \int_0^1\frac{1}{a+bx}\mathrm{d}x+\frac{2\ln a}{b}\cdot \arctan \frac{x}{b}\bigg|_{0}^{+\infty}\\
&=\frac{\pi}{b}\ln \frac{a+b}{a}+\frac{\pi \ln a}{b}=\frac{\pi}{b}\ln \left( a+b \right).
\end{align*}
\end{proof}

\begin{proposition}\label{proposition:Possion积分-数分}
证明Possion积分: \(\int_0^{+\infty} e^{-ax^2} \cos(bx) \mathrm{d}x = \frac{1}{2}\sqrt{\frac{\pi}{a}} e^{-\frac{b^2}{4a}}\), 其中 \(a>0\).
\end{proposition}
\begin{remark}
也可利用留数定理证明(具体见\refpro{Complex Analysis-proposition:几个特殊积分的计算-Poisson积分}).
\end{remark}
\begin{proof}
设
\begin{align}
I(b) &= \int_0^{+\infty} e^{-ax^2} \cos(bx) \mathrm{d}x, \label{myIdef_xyz_123}
\end{align}
其中 \(a>0, b\) 为实数.
由于
\begin{align*}
\left| e^{-ax^2} \cos(bx) \right| \leq e^{-ax^2},
\end{align*}
且根据\hyperref[proposition:重要定积分结论-1]{Gauss-Poisson积分}可知 \(\int_0^{+\infty} e^{-ax^2} \mathrm{d}x\) 收敛, 故$I(b)$关于参数$b$在 \(\mathbb{R}\)上收敛. 
对于任意的 \(b \in \mathbb{R}\) 与 \(x \geq 0\), 有
\begin{align*}
\left| \frac{\partial}{\partial b} \left( e^{-ax^2} \cos(bx) \right) \right|=\left| -x e^{-ax^2} \sin(bx) \right| \leq x e^{-ax^2}.
\end{align*}
由于
\begin{align*}
\int_0^{+\infty} x e^{-ax^2} \mathrm{d}x = \lim_{R \to +\infty} \left[ -\frac{1}{2a} e^{-ax^2} \right]_0^{R} = \frac{1}{2a}<+\infty,
\end{align*}
故由\(M\)判别法可知\(\int_0^{+\infty} f_b(x, b) \mathrm{d}x\)关于参数\(b\) 在 \(\mathbb{R}\) 上一致收敛.
因此，根据\hyperref[theorem:含参反常积分的可微性]{含参反常积分的可微性}和分部积分可得
\begin{align*}
I' (b)&=\int_0^{+\infty}{\frac{\partial}{\partial b}\left( e^{-ax^2}\cos(bx) \right) \mathrm{d}x}=\int_0^{+\infty}{-xe^{-ax^2}\sin(bx)\mathrm{d}x.}
\\
&=\left[ \frac{1}{2a}e^{-ax^2}\sin(bx) \right] _{0}^{+\infty}-\int_0^{+\infty}{\frac{1}{2a}e^{-ax^2}b\cos(bx)\mathrm{d}x}
\\
&=-\frac{b}{2a}\int_0^{+\infty}{e^{-ax^2}\cos(bx)\mathrm{d}x}=-\frac{b}{2a}I(b).
\end{align*}
分离变量并积分, 得
\begin{align*}
\frac{\mathrm{d}I}{I} = -\frac{b}{2a} \mathrm{d}b\Longrightarrow 
\ln|I(b)| = -\frac{b^2}{4a} + C_1. 
\end{align*}
由此可得通解 \(I(b) = C e^{-\frac{b^2}{4a}},\)其中$C,C_1$为某一常数.
下面确定常数 \(C\).由\eqref{myIdef_xyz_123}式和\hyperref[proposition:重要定积分结论-1]{Gauss-Poisson积分}可知
\begin{align*}
C=I(0) = \int_0^{+\infty} e^{-ax^2} \mathrm{d}x= \frac{1}{2}\sqrt{\frac{\pi}{a}}..
\end{align*}
故
\begin{align*}
I(b) = \frac{1}{2}\sqrt{\frac{\pi}{a}} e^{-\frac{b^2}{4a}}.
\end{align*}
\end{proof}

\begin{proposition}\label{proposition:高斯积分推论}
设 \(a>0, b>0\)，则
\begin{align*}
\int_{0}^{+\infty} e^{-a^2 x^2 - \frac{b^2}{x^2}} \mathrm{d}x &= \frac{\sqrt{\pi}}{2a} e^{-2ab}. 
\end{align*}
\end{proposition}
\begin{proof}
令
\begin{align*}
F(b) &= \int_{0}^{+\infty} e^{-a^2 x^2 - \frac{b^2}{x^2}} \mathrm{d}x.
\end{align*}
其中 \(a > 0, b > 0\)。对参数 \(b\) 求导时，取任意的有界闭区间 \([\delta, M]\)（其中 \(\delta > 0\)）。对于任意 \(b \in [\delta, M]\)，易知
\begin{align*}
\left| \frac{\partial}{\partial b} e^{-a^2 x^2 - \frac{b^2}{x^2}} \right| = \left| -\frac{2b}{x^2} e^{-a^2 x^2} e^{-\frac{b^2}{x^2}} \right| \leq 2M x^{-2} e^{-a^2 x^2} e^{-\frac{\delta^2}{x^2}}.
\end{align*}
注意到
\begin{align*}
\underset{x\rightarrow 0^+}{\lim}\left( 2Mx^{-2}e^{-a^2x^2}e^{-\frac{\delta ^2}{x^2}} \right) =0,\quad \underset{x\rightarrow +\infty}{\lim}\left( 2Mx^{-2}e^{-a^2x^2}e^{-\frac{\delta ^2}{x^2}} \right) =0,
\end{align*}
因此$\int_0^{+\infty}{2Mx^{-2}e^{-a^2x^2}e^{-\frac{\delta ^2}{x^2}}\mathrm{d}x}$收敛,故根据 \(M\) 判别法可知 \(\int_0^{+\infty} f_b(x, b) \mathrm{d}x\) 关于 \(b\) 在 \([\delta, M]\) 上一致收敛,再由$\delta$的任意性知 \(\int_0^{+\infty} f_b(x, b) \mathrm{d}x\) 关于 \(b\) 在$[0,M]$上内闭一致收敛.
从而由\hyperref[theorem:含参反常积分的可微性]{含参反常积分的可微性}得
\begin{align*}
F' (b)&=\int_0^{+\infty}{\frac{\partial}{\partial b}e^{-a^2x^2-\frac{b^2}{x^2}}\mathrm{d}x}=\int_0^{+\infty}{-\frac{2b}{x^2}e^{-a^2x^2-\frac{b^2}{x^2}}\mathrm{d}x}
\\
&\xlongequal{x=\frac{b}{at}}\int_{+\infty}^0{\left( -\frac{2a^2t^2}{b} \right) e^{-a^2t^2-\frac{b^2}{t^2}}}\cdot \left( -\frac{b}{at^2} \right) \mathrm{d}t,
\\
&=-\int_0^{+\infty}{2ae^{-a^2t^2-\frac{b^2}{t^2}}\mathrm{d}t=}-2aF\left( b \right) .
\end{align*}
解上述一阶线性齐次常微分方程 \(\frac{\mathrm{d}F}{\mathrm{d}b} = -2a F\) 得通解 \(F(b) = C e^{-2ab}\)。
下面确定常数 \(C\).令 \(b \to 0^+\),由\hyperref[theorem:含参量反常积分的连续性]{含参量反常积分的连续性}可得
\begin{align*}
F(0) = \int_{0}^{+\infty} e^{-a^2 x^2} \mathrm{d}x = \frac{\sqrt{\pi}}{2a} = C.
\end{align*}
故
\begin{align*}
F(b)=\int_{0}^{+\infty} e^{-a^2 x^2 - \frac{b^2}{x^2}} \mathrm{d}x &= \frac{\sqrt{\pi}}{2a} e^{-2ab}. 
\end{align*}
\end{proof}

\begin{example}
求含参变量瑕积分 \(\int_0^{+\infty} \frac{\cos \beta x}{x^2 + \alpha^2} \mathrm{d}x\) （\(\alpha > 0, \beta > 0\)）.
\end{example}
\begin{note}
提示 \(\frac{1}{x^2 + \alpha^2} = \int_{0}^{+\infty} e^{-t(x^2+\alpha^2)} \mathrm{d}t\).
\end{note}
\begin{proof}
设 \(I(\beta) = \int_0^{+\infty} \frac{\cos \beta x}{x^2 + \alpha^2} \mathrm{d}x\)。任取 \(\varepsilon > 0\)，定义
\begin{align*}
I(\beta )\triangleq \int_0^{+\infty}{\int_0^{+\infty}{e^{-t\left( x^2+\alpha ^2 \right)}\cos \left( \beta x \right) \mathrm{d}t\mathrm{d}x}}.
\end{align*}
现在验证$I(\beta)$满足\cref{theorem:含参积分可积性定理}的条件.令 
\begin{align*}
f(x, t) = e^{-t(x^2+\alpha^2)} \cos(\beta x),\quad (x,t)\in [0, +\infty) \times [\varepsilon, +\infty).
\end{align*}
由\hyperref[proposition:重要定积分结论-1]{Gauss-Poisson积分}知,对$\forall \delta>0$,当$t\in[\delta,+\infty)$时,有
\begin{align*}
\int_0^{+\infty}{\left| f\left( x,t \right) \right|\mathrm{d}x}\leqslant \int_0^{+\infty}{e^{-\delta \left( x^2+\alpha ^2 \right)}\mathrm{d}x}\xlongequal{u=\sqrt{\delta}x}\frac{e^{-\delta \alpha ^2}}{\sqrt{\delta}}\int_0^{+\infty}{e^{-u^2}\mathrm{d}u}=\frac{e^{-\delta \alpha ^2}}{2}\sqrt{\frac{\pi}{\delta}},
\end{align*}
故由$M$判别法知\(\int_0^{+\infty} f(x,t) \mathrm{d}x\) 在 \([\delta, +\infty)\) 上关于 \(t\)一致收敛,再由$\delta$的任意性知\(\int_0^{+\infty} f(x,t) \mathrm{d}x\) 在 \([0, +\infty)\) 上关于 \(t\)内闭一致收敛.
又当$x\in[0,+\infty)$时,有
\begin{align*}
\int_0^{+\infty}{\left| f\left( x,t \right) \right|\mathrm{d}t}\leqslant \int_0^{+\infty}{e^{-t\left( x^2+\alpha ^2 \right)}\mathrm{d}t}=\frac{1}{x^2+\alpha ^2}\leqslant \frac{1}{\alpha ^2},
\end{align*}
故由$M$判别法知\(\int_0^{+\infty}{f\left( x,t \right) \mathrm{d}t}\) 在 \([0, +\infty)\) 上关于 \(x\)内闭一致收敛。
并且
\begin{align*}
\int_0^{+\infty}{\mathrm{d}x\int_0^{+\infty}{|f(x,t)|\mathrm{d}t}}\leqslant \int_0^{+\infty}{\frac{1}{x^2+\alpha ^2}\mathrm{d}x}<+\infty .
\end{align*}
综上所述，根据\cref{theorem:含参积分可积性定理}可得
\begin{align*}
I(\beta) = \int_{0}^{+\infty} e^{-t\alpha^2} \left( \int_{0}^{+\infty} e^{-t x^2} \cos(\beta x) \mathrm{d}x \right) \mathrm{d}t.
\end{align*}
再利用\cref{proposition:Possion积分-数分}可得
\begin{align*}
I(\beta) = \frac{\sqrt{\pi}}{2} \int_{0}^{+\infty} \frac{1}{\sqrt{t}} e^{-\alpha^2 t - \frac{\beta^2}{4t}} \mathrm{d}t \xlongequal{t=u^2}\sqrt{\pi}\int_0^{+\infty}{e^{-\alpha ^2u^2-\frac{\beta ^2}{4u^2}}\mathrm{d}u}.
\end{align*}
最后，由\cref{proposition:高斯积分推论}可得
\begin{align*}
I(\beta) = \sqrt{\pi} \cdot \frac{\sqrt{\pi}}{2\alpha} e^{-2\alpha \cdot \frac{\beta}{2}}
= \frac{\pi}{2\alpha} e^{-\alpha\beta}. 
\end{align*}
\end{proof}

\begin{example}
设 $|a| \neq 1$，计算
\begin{align*}
\int_{0}^{\frac{\pi}{2}} \frac{\arctan(a \tan x)}{\tan x} \, \mathrm{d}x.
\end{align*}
\end{example}
\begin{solution}
令
\begin{align*}
I(a)=\int_{0}^{\frac{\pi}{2}}\frac{\arctan(a\tan x)}{\tan x}\,\mathrm{d}x.
\end{align*}
显然 $I(0)=0$,且$I(a)$在$\mathbb{R}\setminus \{-1,1\}$上收敛,被积函数关于$a$求偏导后的积分在$[0,\frac{\pi}{2}]$上一致收敛.于是由\hyperref[theorem:含参反常积分的可微性]{含参反常积分的可微性},再记 $c=|a|$，则
\begin{align*}
I' (a)=\int_0^{\frac{\pi}{2}}{\frac{\mathrm{d}x}{1+a^2\tan ^2x}}\xlongequal{t=\tan x}\int_0^{\infty}{\frac{\mathrm{d}t}{(1+a^2t^2)(1+t^2)}}=\frac{1}{1-c^2}\int_0^{\infty}{\left( \frac{1}{1+t^2}-\frac{c^2}{1+c^2t^2} \right) \mathrm{d}t}=\frac{\pi}{2(1+|a|)}.
\end{align*}
因此，若 $a>0$，则
\begin{align*}
I(a)=\int_{0}^{a}\frac{\pi}{2(1+s)}\,\mathrm{d}s
=\frac{\pi}{2}\ln(1+a);
\end{align*}
若 $a<0$，则
\begin{align*}
I(a)=\int_{0}^{a}\frac{\pi}{2(1-s)}\,\mathrm{d}s
=-\frac{\pi}{2}\ln(1-a).
\end{align*}
综上，
\begin{align*}
I(a)=
\begin{cases}
\dfrac{\pi}{2}\ln(1+a), & a>0,\\[6pt]
0, & a=0,\\[6pt]
-\dfrac{\pi}{2}\ln(1-a), & a<0,
\end{cases}
\quad |a|\neq 1.
\end{align*}
也可写成
\begin{align*}
I(a)=\frac{\pi}{2}\operatorname{sgn}(a)\ln(1+|a|),
\end{align*}
其中 $\operatorname{sgn}(0)=0$。
\end{solution}



\subsection{级数展开方法}

\begin{proposition}[级数逐项可积等价于余项趋于0]\label{proposition:级数逐项可积等价于余项趋于0}
设${f_n}_{n=1}^{\infty}$都在$[a,b]$上可积，且$\sum\limits_{n=1}^{\infty} f_n$在$[a,b]$上可积，且$\sum\limits_{n=1}^{\infty}\int_a^b f_n(x)\,\mathrm{d}x$收敛，则
\begin{align*}
\sum_{n=1}^{\infty}\int_a^b f_n(x)\,\mathrm{d}x
= \int_a^b \sum_{n=1}^{\infty} f_n(x)\,\mathrm{d}x\iff \lim_{m\to\infty}\int_a^b \sum_{n=m}^{\infty} f_n(x)\,\mathrm{d}x = 0.
\end{align*}
\end{proposition}
\begin{proof}
记$S_m(x)=\sum\limits_{n=1}^{m} f_n(x)$，$S(x)=\sum\limits_{n=1}^{\infty} f_n(x)$，并设$A=\sum\limits_{n=1}^{\infty}\int_a^b f_n(x)\,\mathrm{d}x$。
由于有限和可积且积分可加，有
\begin{align*}
\sum_{n=1}^{m}\int_a^b f_n(x)\,\mathrm{d}x
= \int_a^b S_m(x)\,\mathrm{d}x.
\end{align*}
因此
\begin{align*}
\sum_{n=1}^{\infty}\int_a^b f_n(x)\,\mathrm{d}x
= \int_a^b S(x)\,\mathrm{d}x
&\Longleftrightarrow
\lim_{m\to\infty}\int_a^b S_m(x)\,\mathrm{d}x
= \int_a^b S(x)\,\mathrm{d}x \\
&\Longleftrightarrow
\lim_{m\to\infty}\int_a^b \bigl(S(x)-S_m(x)\bigr)\,\mathrm{d}x = 0.
\end{align*}
而
\begin{align*}
S(x)-S_m(x)=\sum_{n=m+1}^{\infty} f_n(x),
\end{align*}
故结论成立.
\end{proof}

\begin{example}
计算$\int_{0}^{\infty} \frac{x}{1 + e^{x}} \mathrm{d}x$. 
\end{example}
\begin{solution}
由于
\begin{align*}
\frac{1}{1 + x} = \sum_{n = 0}^{\infty} (-1)^n x^n, \quad 0 < x < 1.
\end{align*}
并且\(0 < e^{-x} < 1\)，故
\begin{align}
\int_0^{+\infty} \frac{x}{1 + e^x} \mathrm{d}x &= \int_0^{+\infty} \frac{x e^{-x}}{1 + e^{-x}} \mathrm{d}x = \int_0^{+\infty} x \sum_{n = 0}^{\infty} (-1)^n e^{-(n + 1)x} \mathrm{d}x  \\
&\xlongequal{\text{\hyperlink{积分求和换序equation100.39}{换序}}} \sum_{n = 0}^{\infty} (-1)^n \int_0^{+\infty} x e^{-(n + 1)x} \mathrm{d}x = \sum_{n = 0}^{\infty} \frac{(-1)^n}{(n + 1)^2}  \\
&= \sum_{n = 0}^{\infty} \frac{1}{(2n + 1)^2} - \sum_{n = 0}^{\infty} \frac{1}{(2n)^2}.\label{equation:100.3911222} 
\end{align}
又因为\(\sum_{n = 1}^{\infty} \frac{1}{n^2} = \frac{\pi^2}{6}\)，所以
\begin{align*}
\sum_{n = 0}^{\infty} \frac{1}{(2n)^2} &= \frac{1}{4} \sum_{n = 1}^{\infty} \frac{1}{n^2} = \frac{\pi^2}{24}, \\
\sum_{n = 0}^{\infty} \frac{1}{(2n + 1)^2} &= \sum_{n = 1}^{\infty} \frac{1}{n^2} - \sum_{n = 0}^{\infty} \frac{1}{(2n)^2} = \frac{\pi^2}{8}.
\end{align*}
故
\begin{align*}
\int_0^{+\infty} \frac{x}{1 + e^x} \mathrm{d}x &= \sum_{n = 0}^{\infty} \frac{1}{(2n + 1)^2} - \sum_{n = 0}^{\infty} \frac{1}{(2n)^2} = \frac{\pi^2}{8} - \frac{\pi^2}{24} = \frac{\pi^2}{12}.
\end{align*}
\hypertarget{积分求和换序equation100.39}{下面证明}\eqref{equation:100.3911222}式换序成立，由\cref{proposition:级数逐项可积等价于余项趋于0}知这等价于证明\(\lim_{m \to +\infty} \int_0^{+\infty} \sum_{n = m}^{\infty} x (-1)^n e^{-(n + 1)x} \mathrm{d}x = 0\).
由\hyperref[theorem:交错级数不等式]{交错级数不等式}及\(x e^{-(n + 1)x}\)关于\(n\)非负递减，对\(\forall m \in \mathbb{N}\)，都有
\\
\begin{align*}
\int_0^{+\infty} \left| \sum_{n = m}^{\infty} x (-1)^n e^{-(n + 1)x} \right| \mathrm{d}x &\leqslant \int_0^{+\infty} x e^{-(m + 1)x} \mathrm{d}x = -\frac{x e^{-(m + 1)x}}{m + 1} \Big|_0^{+\infty} + \frac{1}{m + 1} \int_0^{+\infty} e^{-(m + 1)x} \mathrm{d}x = \frac{1}{(m + 1)^2}.
\end{align*}
令\(m \to +\infty\)，得\(\lim_{m \to +\infty} \int_0^{+\infty} \sum_{n = m}^{\infty} x (-1)^n e^{-(n + 1)x} \mathrm{d}x = 0\). 故\eqref{equation:100.3911222}式换序成立. 
\end{solution}

\begin{example}
证明:
\begin{align*}
\int_{0}^{+\infty} \frac{x}{e^x - e^{-x}} \, \mathrm{d}x = \frac{\pi^2}{8}.
\end{align*}
\end{example}
\begin{proof}
由Taylor定理知
\begin{align}
\int_0^{+\infty}{\frac{x}{e^x-e^{-x}}\mathrm{d}x}
= \int_0^{+\infty}{\frac{xe^{-x}}{1-e^{-2x}}\mathrm{d}x}
= \int_0^{+\infty}{\sum_{n=0}^{\infty}{xe^{-(2n+1)x}}\,\mathrm{d}x}.
\label{eq:32_0}
\end{align}
现在证明
\begin{align}
\int_0^{+\infty}{\sum_{n=0}^{\infty}{xe^{-(2n+1)x}}\,\mathrm{d}x}
= \sum_{n=0}^{\infty}{\int_0^{+\infty}{xe^{-(2n+1)x}\,\mathrm{d}x}}.
\label{eq:32_1}
\end{align}
由\cref{proposition:级数逐项可积等价于余项趋于0}知只需证明
\begin{align*}
\lim_{m\rightarrow +\infty}\int_0^{+\infty}{\sum_{n=m}^{\infty}{xe^{-(2n+1)x}}\,\mathrm{d}x}=0.
\end{align*}
由L'Hospital法则知$\lim\limits_{x\rightarrow 0^+}\frac{x}{1-e^{-2x}}=\frac{1}{2}$，故$\exists \delta >0$，使得
\begin{align*}
\frac{x}{1-e^{-2x}}\leqslant 1,\quad \forall x\in[0,\delta].
\end{align*}
进而再利用等比求和公式得，对$\forall m\in\mathbb{N}$，有
\begin{align*}
\int_0^{+\infty}{\sum_{n=m}^{\infty}{xe^{-(2n+1)x}}\,\mathrm{d}x}
&= \int_0^{+\infty}{\frac{xe^{-(2m+1)x}}{1-e^{-2x}}\,\mathrm{d}x} = \int_0^{\delta}{\frac{xe^{-(2m+1)x}}{1-e^{-2x}}\,\mathrm{d}x}
+\int_{\delta}^{+\infty}{\frac{xe^{-(2m+1)x}}{1-e^{-2x}}\,\mathrm{d}x} \\
&\leqslant \int_0^{\delta}{e^{-(2m+1)x}\,\mathrm{d}x}
+\frac{1}{1-e^{-2\delta}}\int_{\delta}^{+\infty}{xe^{-(2m+1)x}\,\mathrm{d}x} \\
&\xlongequal{\text{分部积分}}
\frac{1-e^{-(2m+1)\delta}}{2m+1}
+\frac{1}{1-e^{-2\delta}}\cdot \frac{\delta e^{-(2m+1)\delta}}{2m+1}
+\frac{1}{1-e^{-2\delta}}\cdot \frac{e^{-(2m+1)\delta}}{(2m+1)^2}.
\end{align*}
令$m\rightarrow +\infty$，得
\begin{align*}
\lim_{m\rightarrow +\infty}\int_0^{+\infty}{\sum_{n=m}^{\infty}{xe^{-(2n+1)x}}\,\mathrm{d}x}=0.
\end{align*}
故\eqref{eq:32_1}式成立。因此再结合\eqref{eq:32_0}式可得
\begin{align*}
\int_0^{+\infty}{\frac{x}{e^x-e^{-x}}\mathrm{d}x}
&= \int_0^{+\infty}{\sum_{n=0}^{\infty}{xe^{-(2n+1)x}}\,\mathrm{d}x} = \sum_{n=0}^{\infty}{\int_0^{+\infty}{xe^{-(2n+1)x}\,\mathrm{d}x}} \\
&\xlongequal{\text{分部积分}}
\sum_{n=0}^{\infty}{\int_0^{+\infty}{\frac{e^{-(2n+1)x}}{2n+1}\,\mathrm{d}x}}
= \sum_{n=0}^{\infty}{\frac{1}{(2n+1)^2}} \\
&= \sum_{n=1}^{\infty}{\frac{1}{n^2}}-\sum_{n=1}^{\infty}{\frac{1}{(2n)^2}}
= \sum_{n=1}^{\infty}{\frac{1}{n^2}}-\frac{1}{4}\sum_{n=1}^{\infty}{\frac{1}{n^2}} \\
&= \frac{3}{4}\sum_{n=1}^{\infty}{\frac{1}{n^2}}
= \frac{\pi^2}{8}.
\end{align*}
\end{proof}

\begin{example}
\begin{align*}
\int_{0}^{+\infty}{xe^{-\frac{x^2}{2}}\left( \sum_{n=0}^{\infty}{\frac{x^{2n}}{\left[ (2n)!! \right]^2}} \right) \,\mathrm{d}x}=\sqrt{e}.
\end{align*}
\end{example}
\begin{proof}
首先，反复分部积分可得
\begin{align}
\int_0^{+\infty}{t^ne^{-t}\,\mathrm{d}t}
&\xlongequal{\text{分部积分}} t^ne^{-t}\big|_{+\infty}^{0}+n\int_0^{+\infty}{t^{n-1}e^{-t}\,\mathrm{d}t}= n\int_0^{+\infty}{t^{n-1}e^{-t}\,\mathrm{d}t} \nonumber \\
&\xlongequal{\text{分部积分}} n(n-1)\int_0^{+\infty}{t^{n-2}e^{-t}\,\mathrm{d}t} = \cdots = n!.
\label{eq:41_0}
\end{align}
现在证明
\begin{align}
\int_0^{+\infty}{\sum_{n=0}^{\infty}{\frac{t^ne^{-t}}{2^n(n!)^2}}\,\mathrm{d}t}
= \sum_{n=0}^{\infty}{\frac{1}{2^n(n!)^2}\int_0^{+\infty}{t^ne^{-t}\,\mathrm{d}t}}.
\label{eq:41_1}
\end{align}
由\cref{proposition:级数逐项可积等价于余项趋于0}知只需证明
\begin{align*}
\lim_{m\rightarrow +\infty}\int_0^{+\infty}{\sum_{n=m}^{\infty}{\frac{t^ne^{-t}}{2^n(n!)^2}}\,\mathrm{d}t}=0.
\end{align*}
对$\forall \varepsilon >0$，取$A>0$，满足$e^{-\frac{A}{2}}<\varepsilon$. 从而由Taylor公式可得，对$\forall m\in \mathbb{N}$，有
\begin{align}
\int_A^{+\infty}{\sum_{n=m}^{\infty}{\frac{t^ne^{-t}}{2^n(n!)^2}}\,\mathrm{d}t}
&\leqslant \int_A^{+\infty}{\sum_{n=0}^{\infty}{\frac{t^ne^{-t}}{2^n(n!)^2}}\,\mathrm{d}t} \leqslant \int_A^{+\infty}{\sum_{n=0}^{\infty}{\frac{t^ne^{-t}}{2^nn!}}\,\mathrm{d}t} \nonumber \\
&= \int_A^{+\infty}{e^{-t}\sum_{n=0}^{\infty}{\frac{\left( \frac{t}{2} \right)^n}{n!}}\,\mathrm{d}t} \xlongequal{\text{Taylor公式}} \int_A^{+\infty}{e^{-t}e^{\frac{t}{2}}\,\mathrm{d}t} \nonumber \\
&= \int_A^{+\infty}{e^{-\frac{t}{2}}\,\mathrm{d}t}=e^{-\frac{A}{2}}<\varepsilon.
\label{eq:32_11}
\end{align}
注意到
\begin{align*}
\left| \frac{t^ne^{-t}}{2^n(n!)^2} \right|\leqslant \frac{A^n}{2^n(n!)^2},\quad \forall t\in[0,A].
\end{align*}
显然$\sum\limits_{n=0}^{\infty}{\frac{A^n}{2^n(n!)^2}}$收敛，故由\hyperref[theorem:Weierstrass(魏尔斯特拉斯)判别法]{M判别法}知$\sum\limits_{n=0}^{\infty}{\frac{t^ne^{-t}}{2^n(n!)^2}}$在$[0,A]$上一致收敛. 因此存在$N\in \mathbb{N}$，使得对$\forall m>N$，有
\begin{align*}
\sum_{n=m}^{\infty}{\frac{t^ne^{-t}}{2^n(n!)^2}}<\varepsilon
\Longrightarrow \int_0^A{\sum_{n=m}^{\infty}{\frac{t^ne^{-t}}{2^n(n!)^2}}\,\mathrm{d}t}<\varepsilon A.
\label{eq:32_222222}
\end{align*}
从而结合\eqref{eq:32_11} \eqref{eq:32_222222}式知，对$\forall m>N$，都有
\begin{align*}
\int_0^{+\infty}{\sum_{n=m}^{\infty}{\frac{t^ne^{-t}}{2^n(n!)^2}}\,\mathrm{d}t}
= \int_0^A{\sum_{n=N}^{\infty}{\frac{t^ne^{-t}}{2^n(n!)^2}}\,\mathrm{d}t}
+\int_A^{+\infty}{\sum_{n=m}^{\infty}{\frac{t^ne^{-t}}{2^n(n!)^2}}\,\mathrm{d}t}
<\varepsilon A+\varepsilon.
\end{align*}
即$\lim\limits_{m\rightarrow +\infty}\int_0^{+\infty}{\sum_{n=m}^{\infty}{\frac{t^ne^{-t}}{2^n(n!)^2}}\,\mathrm{d}t}=0$，故\eqref{eq:41_1}式成立. 于是由\eqref{eq:41_0}\eqref{eq:41_1}式和Taylor公式可得
\begin{align*}
\int_0^{+\infty}{xe^{-\frac{x^2}{2}}\left( \sum_{n=0}^{\infty}{\frac{x^{2n}}{[(2n)!!]^2}} \right) \,\mathrm{d}x}
&\xlongequal{t=\frac{x^2}{2}}
\int_0^{+\infty}{e^{-t}\left( \sum_{n=0}^{\infty}{\frac{2^nt^n}{(2^n\cdot n!)^2}} \right) \,\mathrm{d}t} 
= \int_0^{+\infty}{\sum_{n=0}^{\infty}{\frac{t^ne^{-t}}{2^n(n!)^2}}\,\mathrm{d}t} \\
&= \sum_{n=0}^{\infty}{\frac{1}{2^n(n!)^2}\int_0^{+\infty}{t^ne^{-t}\,\mathrm{d}t}} = \sum_{n=0}^{\infty}{\frac{1}{2^nn!}}
\xlongequal{\text{Taylor公式}} e^{\frac{1}{2}}.
\end{align*}
\end{proof}

\begin{proposition}\label{proposition:常用级数公式1}
证明:
\begin{enumerate}[(1)]
\item \(\sum_{n = 1}^{\infty} \frac{q^{n} \sin(nx)}{n}= \arctan \frac{q \sin x}{1 - q \cos x}, |q| \leqslant 1\).

\item \(\sum_{n = 1}^{\infty} \frac{q^{n} \cos(nx)}{n}= -\frac{1}{2} \ln(1 + q^{2} - 2q \cos x), |q| \leqslant 1\).

\item \(\sum_{n = 1}^{\infty} \frac{q^{n} \cos(nx)}{n!}= e^{q \cos x} \cos(q \sin x) - 1, |q| \leqslant 1, x \in \mathbb{R}\).

\item \(\sum_{n = 1}^{\infty} \frac{q^{n} \sin(nx)}{n!}= e^{q \cos x} \sin(q \sin x), |q| \leqslant 1, x \in \mathbb{R}\).
\end{enumerate}
\end{proposition}
\begin{note}
在\(\mathbb{C}\)上，
\begin{align*}
\text{Ln}z &= \ln|z| + i(\text{arg }z + 2k\pi), k \in \mathbb{Z}.  
\end{align*}
我们定义主值支
\begin{align*}
\ln z &= \ln|z| + i\text{arg }z. 
\end{align*}
本部分内容无需记忆，只需要大概有个可以算的感觉即可，实际做题中可以围绕这种级数给出构造.
\end{note}
\begin{proof}
$\Im$表示取虚部,$\Re$表示取实部.
\begin{enumerate}[(1)]
\item 利用欧拉公式有
\begin{align*}
\sum_{n = 1}^{\infty} \frac{q^{n} \sin(nx)}{n} &= \Im\left( \sum_{n = 1}^{\infty} \frac{q^{n} e^{inx}}{n} \right) = \Im\left( \sum_{n = 1}^{\infty} \frac{(q e^{ix})^{n}}{n} \right) = \Im(-\ln(1 - q e^{ix})) \\
&= -\Im\left( \ln|1 - q e^{ix}| + i \frac{-q \sin x}{1 - q \cos x} \right) = \arctan \frac{q \sin x}{1 - q \cos x}.
\end{align*}

\item 利用欧拉公式有
\begin{align*}
\sum_{n = 1}^{\infty} \frac{q^{n} \cos(nx)}{n} &= -\Re\left( \ln|1 - q e^{ix}| + i \frac{-q \sin x}{1 - q \cos x} \right) = -\frac{1}{2} \ln\left[(1 - q \cos x)^{2} + q^{2} \sin^{2} x\right] \\
&= -\frac{1}{2} \ln(1 + q^{2} - 2q \cos x).
\end{align*}

\item 利用欧拉公式有
\begin{align*}
\sum_{n = 1}^{\infty} \frac{q^{n} \cos(nx)}{n!} &= \Re\left( \sum_{n = 1}^{\infty} \frac{(q e^{ix})^{n}}{n!} \right) = \Re\left( e^{q e^{ix}} - 1 \right) =\Re \left( e^{q\cos x+\mathrm{i}q\sin x}-1 \right) 
\\
&= \Re\left( e^{q \cos x} \cos(q \sin x) - 1 + i e^{q \cos x} \sin(q \sin x) \right) \\
&= e^{q \cos x} \cos(q \sin x) - 1.
\end{align*}

\item 利用(3)有
\begin{align*}
\sum_{n = 1}^{\infty} \frac{q^{n} \sin(nx)}{n!} = \Im\left( e^{q \cos x} \cos(q \sin x) - 1 + i e^{q \cos x} \sin(q \sin x) \right) 
= e^{q \cos x} \sin(q \sin x).
\end{align*} 
\end{enumerate}
\end{proof}

\begin{example}
计算
\begin{enumerate}
\item \(\int_{0}^{2\pi} e^{\cos x} \cos(\sin x) \mathrm{d}x\).

\item \(\int_{0}^{\pi} \ln(1 - 2a\cos x + a^{2}) \mathrm{d}x, a \in (0, +\infty) \setminus \{1\}\). 
\end{enumerate}
\end{example}
\begin{remark}
由1的证明可得
\begin{align*}
e^{\cos x}\cos(\sin x) &= \mathrm{Re}\left( \sum_{n = 0}^{\infty} \frac{(e^{\mathrm{i}x})^n}{n!} \right) = \mathrm{Re}\left( \sum_{n = 0}^{\infty} \frac{e^{\mathrm{i}nx}}{n!} \right) = \sum_{n = 0}^{\infty} \frac{\cos(nx)}{n!}.
\end{align*}
实际上，上式就是\nrefpro{proposition:常用级数公式1}{(3)}的结论. 
\end{remark}
\begin{remark}
第2问也可以用含参积分求导的方法进行计算(这个方法更容易想到).
\end{remark}
\begin{proof}
\begin{enumerate}
\item \begin{align*}
\int_0^{2\pi}{e^{\cos x}\cos \left( \sin x \right) \mathrm{d}x}&=\mathrm{Re}\left( \int_0^{2\pi}{e^{\cos x}e^{\mathrm{i}\sin x}\mathrm{d}x} \right) =\mathrm{Re}\left( \int_0^{2\pi}{e^{\cos x+\mathrm{i}\sin x}\mathrm{d}x} \right) 
\\
&=\mathrm{Re}\left( \int_0^{2\pi}{e^{e^{\mathrm{i}x}}\mathrm{d}x} \right) =\mathrm{Re}\left[ \int_0^{2\pi}{\sum_{n=0}^{+\infty}{\frac{\left( e^{\mathrm{i}x} \right) ^n}{n!}}\mathrm{d}x} \right] =\mathrm{Re}\left[ \sum_{n=0}^{+\infty}{\int_0^{2\pi}{\frac{\left( e^{\mathrm{i}x} \right) ^n}{n!}\mathrm{d}x}} \right] 
\\
&=\mathrm{Re}\left( \sum_{n=0}^{+\infty}{\int_0^{2\pi}{\frac{e^{\mathrm{i}nx}}{n!}\mathrm{d}x}} \right) =\mathrm{Re}\left( \int_0^{2\pi}{\frac{e^{\mathrm{i}\cdot 0\cdot x}}{n!}\mathrm{d}x}+\sum_{n=1}^{+\infty}{\frac{e^{2\pi \mathrm{i}x}-1}{\mathrm{i}n\cdot n!}} \right) 
\\
&=\mathrm{Re}\left( \int_0^{2\pi}{1\mathrm{d}x}+0 \right) =2\pi .
\end{align*}

\item 注意到当\(a\in(0,1)\)时，有
\begin{align*}
\sum_{n = 1}^{\infty} \frac{a^n\cos(nx)}{n} &= \mathrm{Re}\left[ \sum_{n = 1}^{\infty} \frac{(ae^{\mathrm{i}x})^n}{n} \right] = -\mathrm{Re}\left[ \ln(1 - ae^{\mathrm{i}x}) \right] \\
&= -\mathrm{Re}\left[ \ln|1 - ae^{\mathrm{i}x}| + \mathrm{i} \arg(1 - ae^{\mathrm{i}x}) \right] = -\ln|1 - ae^{\mathrm{i}x}| \\
&= -\ln|(1 - a\cos x) + a\mathrm{i}\sin x| = -\frac{1}{2}\ln(1 + a^2 - 2a\cos x).
\end{align*}
于是当\(a\in(0,1)\)时，就有
\begin{align*}
\int_0^{\pi} \ln(1 - 2a\cos x + a^2) \mathrm{d}x &= -\frac{1}{2}\int_0^{\pi} \sum_{n = 1}^{\infty} \frac{a^n\cos(nx)}{n} \mathrm{d}x = 0.
\end{align*}
若\(a > 1\)，则\(\frac{1}{a}\in(0,1)\)，从而此时我们有
\begin{align*}
\int_0^{\pi} \ln(1 - 2a\cos x + a^2) \mathrm{d}x &= \pi \ln a^2 + \int_0^{\pi} \ln\left( \frac{1}{a^2} - \frac{2}{a}\cos x + 1 \right) \mathrm{d}x = \pi \ln a^2 = 2\pi \ln a.
\end{align*}

又由\(\ln(1 - 2a\cos x + a^2)\)关于\(a\)的偏导存在可知\(\int_0^{\pi} \ln(1 - 2a\cos x + a^2) \mathrm{d}x\)关于\(a\)连续. 于是由
\\
\begin{align*}
\int_0^{\pi} \ln(1 - 2a\cos x + a^2) \mathrm{d}x &= 2\pi \ln a, \quad \forall a > 1.
\end{align*}
可知当\(a = 1\)时，我们有
\\
\begin{align*}
\int_0^{\pi} \ln(2 - 2\cos x) \mathrm{d}x &= \lim_{a \to 1^+} \int_0^{\pi} \ln(1 - 2a\cos x + a^2) \mathrm{d}x = \lim_{a \to 1^+} (2\pi \ln a) = 0.
\end{align*} 
\end{enumerate}
\end{proof}

\begin{definition}[多重对数函数-$\mathrm{Li}_2$函数]
定义  
\[
\mathrm{Li}_2(x) \triangleq \sum_{n=1}^{\infty} \frac{x^n}{n^2}, \quad x \in [-1,1].
\] 
\end{definition}

\begin{proposition}[$\mathrm{Li}_2$函数的基本性质]\label{proposition:Li_2函数的性质}
\begin{enumerate}[(1)]
\item $\mathrm{Li}_2(x) + \mathrm{Li}_2(1-x) = \frac{\pi^2}{6} - \ln x \cdot \ln(1-x)$，$x \in (0,1)$。  

\item $\mathrm{Li}_2\left( 1 \right) =\sum_{n=1}^{\infty}{\frac{1}{n^2}}=\frac{\pi ^2}{6},\quad \mathrm{Li}_2\left( 0 \right) =0,\quad \mathrm{Li}_2\left(\frac{1}{2}\right) = \frac{\pi^2}{12} - \frac{1}{2} \ln^2 \frac{1}{2}$。
\end{enumerate}
\end{proposition}
\begin{proof}
\begin{enumerate}[(1)]
\item 记 $f(x) \triangleq \mathrm{Li}_2(x)$，$F(x) \triangleq f(x) + f(1-x) + \ln x \ln(1-x)$。显然$f(x)$收敛域为$[-1,1]$,故由\hyperref[theorem:函数项级数的逐项可微性]{函数项级数的逐项可微性}知
\begin{align*}
f'(x) = \sum_{n=1}^{\infty} \frac{x^{n-1}}{n} = -\frac{1}{x} \ln(1-x).
\end{align*}  
于是  
\begin{align*}
F'(x) = -\frac{1}{x} \ln(1-x) + \frac{\ln x}{1-x} - \frac{\ln x}{1-x} + \frac{\ln(1-x)}{x} = 0.
\end{align*}  
故 $F(x) \equiv F(1) = f(0) + f(1) = 0 + \sum_{n=1}^{\infty} \frac{1}{n^2} = \frac{\pi^2}{6}$。  

\item 显然$\mathrm{Li}_2\left( 1 \right) =\sum_{n=1}^{\infty}{\frac{1}{n^2}}=\frac{\pi ^2}{6},\mathrm{Li}_2\left( 0 \right) =0.$由(1)可得  
\begin{align*}
\mathrm{Li}_2\left(\frac{1}{2}\right) + \mathrm{Li}_2\left(\frac{1}{2}\right) = 2\mathrm{Li}_2\left(\frac{1}{2}\right) = \frac{\pi^2}{6} - \ln^2 \frac{1}{2} \implies \mathrm{Li}_2\left(\frac{1}{2}\right) = \frac{\pi^2}{12} - \frac{1}{2} \ln^2 \frac{1}{2}.
\end{align*}  
\end{enumerate}
\end{proof}

\begin{example}
\begin{align*}
\int_{0}^{1} \ln x \ln(1-x) \, \mathrm{d}x = 2 - \frac{\pi^2}{6}.
\end{align*}
\end{example}
\begin{solution}

{\color{blue}解法一:} 由Taylor定理知
\begin{align*}
\ln(1-x) = -\sum\limits_{n=1}^{\infty}\frac{x^n}{n},\quad \forall x\in(0,1).
\end{align*}
由分部积分知
\begin{align*}
\int_0^1 x^n\ln x\,\mathrm{d}x
= \frac{x^{n+1}\ln x}{n+1}\Big|_0^1 - \frac{1}{n+1}\int_0^1 x^n\,\mathrm{d}x
= -\frac{1}{(n+1)^2}.
\end{align*}
从而
\begin{align*}
\int_0^1 \ln x\ln(1-x)\,\mathrm{d}x
&= -\int_0^1 \ln x\left(\sum\limits_{n=1}^{\infty}\frac{x^n}{n}\right)\mathrm{d}x \\
&= -\sum\limits_{n=1}^{\infty}\frac{1}{n}\int_0^1 x^n\ln x\,\mathrm{d}x \\
&= \sum\limits_{n=1}^{\infty}\frac{1}{n(n+1)^2}
= \sum\limits_{n=1}^{\infty}\left[\frac{1}{n(n+1)}-\frac{1}{(n+1)^2}\right] \\
&= \sum\limits_{n=1}^{\infty}\left(\frac{1}{n}-\frac{1}{n+1}\right)+1-\sum\limits_{n=1}^{\infty}\frac{1}{n^2}
= 2-\frac{\pi^2}{6}.
\end{align*}

{\color{blue}解法二:} 由\hyperref[proposition:Li_2函数的性质]{$\mathrm{Li}_2$函数的性质}知
\begin{align*}i
\int_0^1 \ln x\ln(1-x)\,\mathrm{d}x
&= \int_0^1 \left[\frac{\pi^2}{6}-\left(\mathrm{Li}_2(x)+\mathrm{Li}_2(1-x)\right)\right]\mathrm{d}x = \frac{\pi^2}{6}-2\int_0^1 \mathrm{Li}_2(x)\,\mathrm{d}x \\
&= \frac{\pi^2}{6}-2\int_0^1 \sum\limits_{n=1}^{\infty}\frac{x^n}{n^2}\,\mathrm{d}x = \frac{\pi^2}{6}-2\sum\limits_{n=1}^{\infty}\int_0^1 \frac{x^n}{n^2}\,\mathrm{d}x \\
&= \frac{\pi^2}{6}-2\sum\limits_{n=1}^{\infty}\frac{1}{n^2(n+1)} = \frac{\pi^2}{6}-2\sum\limits_{n=1}^{\infty}\left[\frac{1}{n^2}-\frac{1}{n(n+1)}\right] \\
&= \frac{\pi^2}{6}-2\sum\limits_{n=1}^{\infty}\frac{1}{n^2}
+2\sum\limits_{n=1}^{\infty}\left(\frac{1}{n}-\frac{1}{n+1}\right)
= 2-\frac{\pi^2}{6}.
\end{align*}
\end{solution}

\begin{example}
计算$\int_0^{\frac{1}{2}} \frac{\ln x}{1-x}\, \mathrm{d}x$.
\end{example}
\begin{solution}
利用$\ln (1-x)$在$[-1,1]$上的Taylor公式得
\begin{align*}
\int_0^{\frac{1}{2}} \frac{\ln x}{1-x} \, \mathrm{d}x &= \int_{\frac{1}{2}}^1 \frac{\ln(1-x)}{x} \, \mathrm{d}x = -\sum_{n=1}^{\infty} \frac{1}{n} \int_{\frac{1}{2}}^1 x^{n-1} \, \mathrm{d}x\\
&= -\sum_{n=1}^{\infty} \frac{1}{n^2} + \sum_{n=1}^{\infty} \frac{1}{2^n n^2} = -\frac{\pi^2}{6} + \mathrm{Li}_2\left(\frac{1}{2}\right)
\\
&\xlongequal{\text{\hyperref[proposition:Li_2函数的性质]{$\mathrm{Li}_2$函数的基本性质}}}-\frac{\pi^2}{12} - \frac{1}{2} \ln^2 \frac{1}{2}.
\end{align*}

\end{solution}


\subsection{重积分计算}

\begin{theorem}[二重积分换序]\label{theorem:二重积分换序}
证明:\[
\int_{a}^{b} \mathrm{d}x \int_{a}^{x} f(x,y) \mathrm{d}y = \int_{a}^{b} \mathrm{d}y \int_{y}^{b} f(x,y) \mathrm{d}x, \tag{10}
\]
其中\(f(x,y)\)是在由直线\(y = a\), \(x = b\), \(y = x\)所围成的三角形\((\Delta)\)上连续的任意函数.
\end{theorem}
\begin{proof}

\end{proof}

\begin{proposition}\label{proposition:多重积分递推公式}
设$f(x)$在$[a,b]$上连续, 试证: 对任意$x \in (a,b)$, 有
\[
\int_{a}^{x} \mathrm{d}x_1 \int_{a}^{x_1} \mathrm{d}x_2 \cdots \int_{a}^{x_n} f(x_{n + 1}) \mathrm{d}x_{n + 1} = \frac{1}{n!} \int_{a}^{x} (x - y)^n f(y) \mathrm{d}y, \quad n = 1,2,\cdots.
\]
\end{proposition}
\begin{proof}
当$n=1$时,由\hyperref[theorem:二重积分换序]{二重积分换序}可知
\begin{align*}
\int_a^x \mathrm{d}x_1 \int_a^{x_1} f(x_2) \mathrm{d}x_2 = \int_a^x \mathrm{d}x_2 \int_{x_2}^x f(x_2) \mathrm{d}x_1 = \int_a^x (x-x_2) f(x_2) \mathrm{d}x_2 = \int_a^x (x-y) f(y) \mathrm{d}y.
\end{align*}
设原结论对$n=k-1$的情形成立,考虑$n=k$的情形,由归纳假设可知
\begin{align*}
\int_a^{x_1} \mathrm{d}x_1 \cdots \int_a^{x_k} f(x_{k+1}) \mathrm{d}x_{k+1} = \frac{1}{(k-1)!} \int_a^{x_1} (x_1-y)^{k-1} f(y) \mathrm{d}y.
\end{align*}
于是再利用\hyperref[theorem:二重积分换序]{二重积分换序}得
\begin{align*}
\int_a^x \mathrm{d}x_1 \int_a^{x_1} \mathrm{d}x_2 \cdots \int_a^{x_k} f(x_{k+1}) \mathrm{d}x_{k+1} &= \frac{1}{(k-1)!} \int_a^x \mathrm{d}x_1 \int_a^{x_1} (x_1-y)^{k-1} f(y) \mathrm{d}y \\
&= \frac{1}{(k-1)!} \int_a^x \mathrm{d}y \int_y^x (x_1-y)^{k-1} f(y) \mathrm{d}x_1 \\
&= \frac{1}{k!} \int_a^x (x-y)^k f(y) \mathrm{d}y.
\end{align*}
故由数学归纳法知原结论成立.
\end{proof}

\begin{example}
求定义在星形区域 $D = \{(x,y) \mid x^{\frac{2}{3}} + y^{\frac{2}{3}} \leqslant 1\}$ 上满足 $f(1,0) = 1$ 的正值连续函数 $f$ 使得 $\iint\limits_D \frac{f(x,y)}{f(y,x)} \,\mathrm{d}x\mathrm{d}y$ 达到最小，并求出这个最小值.
\end{example}
\begin{solution}
对积分 $I = \iint\limits_D \frac{f(x,y)}{f(y,x)} \,\mathrm{d}x\mathrm{d}y$ 作变换 $x \to y,\ y \to x$，由 $D$ 的对称性，知 $I = \iint\limits_D \frac{f(y,x)}{f(x,y)} \,\mathrm{d}x\mathrm{d}y$. 因而由均值不等式可得
\begin{align*}
I &= \frac{1}{2} \iint\limits_D \left( \frac{f(x,y)}{f(y,x)} + \frac{f(y,x)}{f(x,y)} \right) \,\mathrm{d}x\mathrm{d}y \geqslant \iint\limits_D 1 \,\mathrm{d}x\mathrm{d}y = \sigma(D),
\end{align*}
这里 $\sigma(D)$ 是 $D$ 的面积.
\[
I - \sigma(D) = \frac{1}{2} \iint\limits_D \left( \sqrt{\frac{f(x,y)}{f(y,x)}} - \sqrt{\frac{f(y,x)}{f(x,y)}} \right)^2 \,\mathrm{d}x\mathrm{d}y \geqslant 0.
\]
$I = \sigma(D)$ 当且仅当 $f(x,y) = f(y,x)$. 故所求函数为所有满足 $f(x,y) = f(y,x)$ 及 $f(1,0) = 1$ 的连续正值函数.

$D$ 的边界的参数方程为
\[
x = \cos^3 \varphi,\quad y = \sin^3 \varphi\ (0 \leqslant \varphi \leqslant 2\pi),
\]
故 $I$ 的最小值为
\begin{align*}
\sigma(D) &= \iint\limits_D 1 \,\mathrm{d}x\mathrm{d}y = 4 \iint\limits_{\substack{0 \leqslant r \leqslant 1 \\ 0 \leqslant \varphi \leqslant \frac{\pi}{2}}} 3r \sin^2 \varphi \cos^2 \varphi \,\mathrm{d}r\mathrm{d}\varphi \\
&= 6 \int_0^{\frac{\pi}{2}} \sin^2 \varphi \cos^2 \varphi \,\mathrm{d}\varphi = \frac{3}{8}\pi.
\end{align*}
所以所求最小值是 $\frac{3}{8}\pi$，且当 $f(x,y) = f(y,x)$ 并满足 $f(1,0) = 1$ 时，取到该最小值.

\end{solution}

\begin{example}
求证: $\iint\limits_{[0,1]^2} (xy)^{xy} \, \mathrm{d}x\mathrm{d}y = \int_0^1 t^t \, \mathrm{d}t$.
\end{example}
\begin{proof}
首先化为累次积分
\begin{align*}
\iint\limits_{[0,1]^2} (xy)^{xy} \, \mathrm{d}x\mathrm{d}y = \int_0^1 \mathrm{d}x \int_0^1 (xy)^{xy} \, \mathrm{d}y = \int_0^1 \mathrm{d}x \int_0^x \frac{t^t}{x} \, \mathrm{d}t = \int_0^1 \frac{f(x)}{x} \, \mathrm{d}x,
\end{align*}
其中 $f(x) = \int_0^x t^t \, \mathrm{d}t$. 由分部积分,
\[
\int_0^1 \frac{f(x)}{x} \, \mathrm{d}x = f(x) \ln x \bigg|_0^1 - \int_0^1 x^x \ln x \, \mathrm{d}x = - \int_0^1 x^x \ln x \, \mathrm{d}x.
\]
因为 $(x^x)' = x^x \ln x + x^x$, 所以
\[
\int_0^1 x^x \ln x \, \mathrm{d}x = \int_0^1 \big( (x^x)' - x^x \big) \, \mathrm{d}x = - \int_0^1 x^x \, \mathrm{d}x.
\]
于是
\[
\iint\limits_{[0,1]^2} (xy)^{xy} \, \mathrm{d}x\mathrm{d}y = \int_0^1 t^t \, \mathrm{d}t.
\]
\end{proof}

\begin{example}
计算二重积分 \( I = \iint\limits_D \text{sgn}(x^2 - y^2 + 2) \, \mathrm{d}x\mathrm{d}y \), 其中$D = \{ (x,y) \mid x^2 + y^2 \leqslant 4 \}.$
\end{example}
\begin{solution}
设 \( D \) 在第一象限部分为 \( D_1 \), 则由对称性
\[
I = 4 \iint\limits_{D_1} \text{sgn}(x^2 - y^2 + 2) \, \mathrm{d}x\mathrm{d}y.
\]
设 \( D_2 \) 是 \( D_1 \) 中使得 \( x^2 - y^2 + 2 < 0 \) 的部分, \( D_3 \) 是 \( D_1 \) 中使得 \( x^2 - y^2 + 2 \geqslant 0 \) 的部分, 则 \( D_1 = D_2 \cup D_3 \). 因此
\begin{align*}
I &= 4 \left[ \iint\limits_{D_3} \mathrm{d}x\mathrm{d}y - \iint\limits_{D_2} \mathrm{d}x\mathrm{d}y \right] = 4 [\sigma(D_3) - \sigma(D_2)]
\\
&= 4 \left[ \frac{1}{4} \cdot \pi \cdot 2^2 - 2\sigma(D_2) \right] = 4\pi - 8\sigma(D_2),
\end{align*}
其中 \( \sigma(D_2), \sigma(D_3) \) 分别表示 \( D_2 \) 和 \( D_3 \) 的面积. 在极坐标 \( x = r\cos\varphi, y = r\sin\varphi \) 之下, \( D_2 \) 为
\[
\left\{ (r,\varphi) \mid \frac{\pi}{3} \leqslant \varphi \leqslant \frac{\pi}{2}, \sqrt{-\frac{2}{\cos 2\varphi}} \leqslant r \leqslant 2 \right\}.
\]
因而
\begin{align*}
\sigma(D_2) &= \iint\limits_{D_2} \mathrm{d}x\mathrm{d}y = \int_{\frac{\pi}{3}}^{\frac{\pi}{2}} \mathrm{d}\varphi \int_{\sqrt{-\frac{2}{\cos 2\varphi}}}^2 r \, \mathrm{d}r
\\
&= \frac{1}{2} \int_{\frac{\pi}{3}}^{\frac{\pi}{2}} \left( 4 + \frac{2}{\cos 2\varphi} \right) \mathrm{d}\varphi = \frac{\pi}{3} + \frac{1}{2} \int_{\frac{2\pi}{3}}^{\pi} \frac{1}{\cos\varphi} \mathrm{d}\varphi
\\
&= \frac{\pi}{3} - \frac{1}{2} \ln(2 + \sqrt{3}),
\end{align*}
故
\[
I = \frac{4\pi}{3} + 4\ln(2 + \sqrt{3}).
\]

\end{solution}

\begin{example}
设 \( D = \{ (x,y) \mid x^2 + y^2 \leqslant 1 \} \). 求 \( I = \iint\limits_D \left| \frac{x + y}{\sqrt{2}} - x^2 - y^2 \right| \mathrm{d}x\mathrm{d}y \).
\end{example}
\begin{solution}
由极坐标变换 \( x = r\cos\varphi \), \( y = r\sin\varphi \), \( 0 \leqslant r \leqslant 1 \), \( 0 \leqslant \varphi \leqslant 2\pi \), 有
\begin{align*}
I &= \iint\limits_{\substack{0 \leqslant r \leqslant 1 \\ 0 \leqslant \varphi \leqslant 2\pi}} \left| \frac{\cos\varphi + \sin\varphi}{\sqrt{2}} - r \right| r^2 \mathrm{d}r\mathrm{d}\varphi
\\
&= \iint\limits_{\substack{0 \leqslant r \leqslant 1 \\ 0 \leqslant \varphi \leqslant 2\pi}} \left| \sin\left( \varphi + \frac{\pi}{4} \right) - r \right| r^2 \mathrm{d}r\mathrm{d}\varphi = \iint\limits_{\substack{0 \leqslant r \leqslant 1 \\ 0 \leqslant \varphi \leqslant 2\pi}} \left| \sin\varphi - r \right| r^2 \mathrm{d}r\mathrm{d}\varphi
\\
&= \iint\limits_{\substack{0 \leqslant r \leqslant 1 \\ 0 \leqslant \varphi \leqslant \pi}} \left| \sin\varphi - r \right| r^2 \mathrm{d}r\mathrm{d}\varphi + \iint\limits_{\substack{0 \leqslant r \leqslant 1 \\ \pi \leqslant \varphi \leqslant 2\pi}} \left| \sin\varphi - r \right| r^2 \mathrm{d}r\mathrm{d}\varphi
\\
&= \iint\limits_{\substack{0 \leqslant r \leqslant 1 \\ 0 \leqslant \varphi \leqslant \pi}} \left| \sin\varphi - r \right| r^2 \mathrm{d}r\mathrm{d}\varphi + \iint\limits_{\substack{0 \leqslant r \leqslant 1 \\ 0 \leqslant \varphi \leqslant \pi}} (\sin\varphi + r) r^2 \mathrm{d}r\mathrm{d}\varphi.
\end{align*}
因此, 有
\begin{align*}
I &= \int_0^\pi \mathrm{d}\varphi \int_0^{\sin\varphi} (\sin\varphi - r) r^2 \mathrm{d}r + \int_0^\pi \mathrm{d}\varphi \int_{\sin\varphi}^1 (r - \sin\varphi) r^2 \mathrm{d}r
\\
&\quad + \int_0^\pi \mathrm{d}\varphi \int_0^{\sin\varphi} (\sin\varphi + r) r^2 \mathrm{d}r + \int_0^\pi \mathrm{d}\varphi \int_{\sin\varphi}^1 (\sin\varphi + r) r^2 \mathrm{d}r
\\
&= \int_0^\pi \mathrm{d}\varphi \int_0^{\sin\varphi} 2\sin\varphi \cdot r^2 \mathrm{d}r + \int_0^\pi \mathrm{d}\varphi \int_{\sin\varphi}^1 2r \cdot r^2 \mathrm{d}r
\\
&= \int_0^\pi \frac{2}{3} \sin^4\varphi \mathrm{d}\varphi + \int_0^\pi \frac{1}{2}(1 - \sin^4\varphi) \mathrm{d}\varphi
\\
&= \frac{1}{6} \int_0^\pi \sin^4\varphi \mathrm{d}\varphi + \frac{\pi}{2} = \frac{1}{6} \cdot \frac{3\pi}{8} + \frac{\pi}{2} = \frac{9}{16}\pi.
\end{align*}

\end{solution}

\begin{example}
设 \( f \) 是定义在正方形 \( S = \{ (x,y) \mid 0 \leqslant x \leqslant 1, 0 \leqslant y \leqslant 1 \} \) 上的四阶连续可微函数, 在 \( S \) 的边界上为零, 并且
\[
\left| \frac{\partial^4 f}{\partial x^2 \partial y^2} \right| \leqslant M.
\]
求证: 
\[
\left| \iint_S f(x,y) \, \mathrm{d}x\mathrm{d}y \right| \leqslant \frac{1}{144} M.
\]
\end{example}
\begin{proof}
考虑函数 \( g(x,y) = x(1 - x)y(1 - y) \). 易知
\[
\frac{\partial^4 g}{\partial x^2 \partial y^2} = 4, \quad \iint_S g(x,y) \, \mathrm{d}x\mathrm{d}y = \frac{1}{36}.
\]
因为 \( f \) 在 \( S \) 的边界上为零, 所以 \( \frac{\partial^2 f}{\partial y^2} \) 在 \( x = 0 \) 和 \( x = 1 \) 时为零. 于是
\begin{align*}
&\iint_S{\frac{\partial ^4f}{\partial x^2\partial y^2}}\cdot g\,\mathrm{d}x\mathrm{d}y=\int_0^1{\mathrm{d}y\int_0^1{\frac{\partial ^4f}{\partial x^2\partial y^2}}}\cdot g\,\mathrm{d}x
\\
&=\int_0^1{\mathrm{d}y\left( \left. \frac{\partial ^3f}{\partial x\partial y^2}\cdot g \right|_{x=0}^{1}-\int_0^1{\frac{\partial ^3f}{\partial x\partial y^2}}\cdot \frac{\partial g}{\partial x}\,\mathrm{d}x \right)}
\\
&=-\int_0^1{\mathrm{d}y\int_0^1{\frac{\partial ^3f}{\partial x\partial y^2}}}\cdot \frac{\partial g}{\partial x}\,\mathrm{d}x
\\
&=-\int_0^1{\mathrm{d}y\left( \left. \frac{\partial ^2f}{\partial y^2}\cdot \frac{\partial g}{\partial x} \right|_{x=0}^{1}-\int_0^1{\frac{\partial ^2f}{\partial y^2}}\cdot \frac{\partial ^2g}{\partial x^2}\,\mathrm{d}x \right)}
\\
&=\int_0^1{\mathrm{d}y\int_0^1{\frac{\partial ^2f}{\partial y^2}}}\cdot \frac{\partial ^2g}{\partial x^2}\,\mathrm{d}x
\\
&=\iint_S{\frac{\partial ^2f}{\partial y^2}}\cdot \frac{\partial ^2g}{\partial x^2}\,\mathrm{d}x\mathrm{d}y.
\end{align*}
同理, 由于 \( \frac{\partial^2 g}{\partial x^2} \) 在 \( y = 0 \) 和 \( y = 1 \) 时为零, 作与上面类似的推导, 可得
\[
\iint_S \frac{\partial^4 g}{\partial x^2 \partial y^2} \cdot f \, \mathrm{d}x\mathrm{d}y = \iint_S \frac{\partial^2 f}{\partial y^2} \cdot \frac{\partial^2 g}{\partial x^2} \, \mathrm{d}x\mathrm{d}y.
\]
因此
\[
\iint_S \frac{\partial^4 f}{\partial x^2 \partial y^2} \cdot g \, \mathrm{d}x\mathrm{d}y = \iint_S \frac{\partial^4 g}{\partial x^2 \partial y^2} \cdot f \, \mathrm{d}x\mathrm{d}y.
\]
从而
\begin{align*}
&\left| \iint_S{f\,\mathrm{d}x\mathrm{d}y} \right|=\frac{1}{4}\left| \iint_S{4f\mathrm{d}x\mathrm{d}y} \right|=\frac{1}{4}\left| \iint_S{\frac{\partial ^4g}{\partial x^2\partial y^2}f\mathrm{d}x\mathrm{d}y} \right|
\\
&=\frac{1}{4}\left| \iint_S{\frac{\partial ^4f}{\partial x^2\partial y^2}\cdot g\,\mathrm{d}x\mathrm{d}y} \right|\leqslant \frac{M}{4}\iint_S{g\,\mathrm{d}x\mathrm{d}y}=\frac{M}{144}.
\end{align*}
\end{proof}

\begin{theorem}[Poincaré(庞加莱)不等式]\label{theorem:Poincaré(庞加莱)不等式}
设 \(\varphi, \psi\) 是 \([a,b]\) 上的连续函数, \(f\) 在区域 \(D = \{(x,y) \mid a \leqslant x \leqslant b, \varphi(x) \leqslant y \leqslant \psi(x)\}\) 上连续可微, 且有 \(f(x, \varphi(x)) = 0\ (x \in [a,b])\). 则存在 \(M > 0\), 使得
\[
\iint_D f^2(x,y) \, \mathrm{d}x\mathrm{d}y \leqslant M \iint_D (f_y'(x,y))^2 \, \mathrm{d}x\mathrm{d}y.
\]
\end{theorem}
\begin{proof}
由 Newton-Leibniz 公式和 Cauchy 不等式可得
\begin{align*}
f^2(x,y)&=[f(x,y)-f(x,\varphi (x))]^2=\left( \int_{\varphi (x)}^y{\frac{\partial f}{\partial t}(x,t)\,\mathrm{d}t} \right) ^2
\\
&\leqslant (y-\varphi (x))\int_{\varphi (x)}^y{\left( \frac{\partial f}{\partial t}(x,t) \right) ^2\,\mathrm{d}t},
\end{align*}
因此
\begin{align*}
&\iint_D{f^2\left( x,y \right) \,\mathrm{d}x\mathrm{d}y}=\int_a^b{\mathrm{d}x}\int_{\varphi \left( x \right)}^{\psi \left( x \right)}{f^2\left( x,y \right) \,\mathrm{d}y}
\\
&\leqslant \int_a^b{\mathrm{d}x}\int_{\varphi \left( x \right)}^{\psi \left( x \right)}{\left( y-\varphi \left( x \right) \right) \mathrm{d}y}\int_{\varphi \left( x \right)}^y{\left( \frac{\partial f}{\partial t}\left( x,t \right) \right) ^2\,\mathrm{d}t}
\\
&=\int_a^b{\mathrm{d}x}\int_{\varphi \left( x \right)}^{\psi \left( x \right)}{\left( \frac{\partial f}{\partial t}\left( x,t \right) \right) ^2\mathrm{d}t}\int_t^{\psi \left( x \right)}{\left( y-\varphi \left( x \right) \right) \mathrm{d}y}
\\
&\leqslant \int_a^b{\mathrm{d}x}\int_{\varphi \left( x \right)}^{\psi \left( x \right)}{\left( \frac{\partial f}{\partial t}\left( x,t \right) \right) ^2\mathrm{d}t}\int_{\varphi \left( x \right)}^{\psi \left( x \right)}{\left( y-\varphi \left( x \right) \right) \mathrm{d}y}
\\
&=\int_a^b{\mathrm{d}x}\int_{\varphi \left( x \right)}^{\psi \left( x \right)}{\frac{1}{2}\left( \psi \left( x \right) -\varphi \left( x \right) \right) ^2\left( \frac{\partial f}{\partial t}\left( x,t \right) \right) ^2\mathrm{d}t}
\\
&\leqslant M\int_a^b{\left( \int_{\varphi \left( x \right)}^{\psi \left( x \right)}{\left( \frac{\partial f}{\partial t}\left( x,t \right) \right) ^2\mathrm{d}t} \right) \mathrm{d}x}
\\
&=M\iint_D{\left( \frac{\partial f}{\partial y}\left( x,y \right) \right) ^2\,\mathrm{d}x\mathrm{d}y},
\end{align*}
这里 \(M\) 是满足 \(M > \max_{a \leqslant x \leqslant b} \frac{1}{2} (\psi(x) - \varphi(x))^2\) 的常数.
\end{proof}

\begin{example}
设 \( a > 0 \), \( \Omega_n(a) : x_1 + x_2 + \cdots + x_n \leqslant a, x_i \geqslant 0\ (i = 1,2,\cdots,n) \). 求积分
\[
I_n(a) = \int \cdots \int_{\Omega_n(a)} x_1 x_2 \cdots x_n \mathrm{d}x_1 \mathrm{d}x_2 \cdots \mathrm{d}x_n.
\]
\end{example}
\begin{solution}
作变换 \( x_i = a t_i \), \( i = 1,2,\cdots,n \), 则
\begin{align*}
I_n(a) = a^{2n} \int \cdots \int_{\Omega_n(1)} t_1 t_2 \cdots t_n \mathrm{d}t_1 \mathrm{d}t_2 \cdots \mathrm{d}t_n = a^{2n} I_n(1).
\end{align*}
再用累次积分, 可得
\begin{align*}
I_n(1) &= \int \cdots \int_{\Omega_n(1)} t_1 t_2 \cdots t_n \mathrm{d}t_1 \mathrm{d}t_2 \cdots \mathrm{d}t_n
\\
&= \int_0^1 t_n \mathrm{d}t_n \int \cdots \int_{\substack{t_1 + t_2 + \cdots + t_{n-1} \leqslant 1 - t_n}} t_1 \cdots t_{n-1} \mathrm{d}t_1 \cdots \mathrm{d}t_{n-1}
\\
&= \int_0^1 t_n I_{n-1}(1 - t_n) \mathrm{d}t_n = \int_0^1 t_n (1 - t_n)^{2(n-1)} I_{n-1}(1) \mathrm{d}t_n.
\end{align*}
因此,
\[
I_n(1) = \frac{1}{2n(2n - 1)} I_{n-1}(1).
\]
注意到 \( I_1(1) = \int_0^1 t \mathrm{d}t = \frac{1}{2} \). 由上面的递推公式, 可得 \( I_n(1) = \frac{1}{(2n)!} \). 故 \( I_n(a) = \frac{a^{2n}}{(2n)!} \).

\end{solution}












\subsection{其他}

\begin{example}\label{example:常用定积分公式}
证明积分$\int_{0}^{\infty}e^{-ax^{2}-\frac{b}{x^{2}}}\mathrm{d}x=\frac{1}{2}\sqrt{\frac{\pi}{a}}e^{-2\sqrt{ab}},a,b>0.$
\end{example}
\begin{proof}
当\(a=1\)时，就有
\begin{align*}
\int_0^{+\infty}{e^{-x^2-\frac{b}{x^2}}\mathrm{d}x}&=e^{-2\sqrt{b}}\int_0^{+\infty}{e^{-\left( x-\frac{\sqrt{b}}{x} \right) ^2}\mathrm{d}x}\xlongequal{y=\frac{\sqrt{b}}{x}}e^{-2\sqrt{b}}\int_0^{+\infty}{\frac{\sqrt{b}}{y^2}e^{-\left( \frac{\sqrt{b}}{y}-y \right) ^2}\mathrm{d}y}\nonumber\\
&=\frac{e^{-2\sqrt{b}}}{2}\int_0^{+\infty}{\left( 1+\frac{\sqrt{b}}{y^2} \right) e^{-\left( y-\frac{\sqrt{b}}{y} \right) ^2}\mathrm{d}y}=\frac{e^{-2\sqrt{b}}}{2}\int_0^{+\infty}{e^{-\left( y-\frac{\sqrt{b}}{y} \right) ^2}\mathrm{d}\left( y-\frac{\sqrt{b}}{y} \right)}\nonumber\\
&=\frac{e^{-2\sqrt{b}}}{2}\int_{-\infty}^{+\infty}{e^{-t^2}\mathrm{d}t}=\frac{\sqrt{\pi}}{2}e^{-2\sqrt{b}}.
\end{align*}
于是对\(\forall a>0\)，就有
\begin{align*}
\int_0^{+\infty}{e^{-ax^2-\frac{b}{x^2}}\mathrm{d}x}=\frac{1}{\sqrt{a}}\int_0^{+\infty}{e^{-x^2-\frac{ab}{x^2}}\mathrm{d}x}=\frac{1}{2}\sqrt{\frac{\pi}{a}}e^{-2\sqrt{ab}}.
\end{align*}
\end{proof}

\begin{example}
计算$\int_{0}^{\infty}\frac{\cos(ax)}{1 + x^{2}}\mathrm{d}x, a\in\mathbb{R}$. 
\end{example}
\begin{remark}
本题可以用复变函数的方法(留数定理)来计算. 但是我们这里用基本的高等数学的方法来计算.

$\int_0^{\infty}{\frac{\sin \left( ax \right)}{1+x^2}\mathrm{d}x}$这个积分没办法算出具体的初等数值.
\end{remark}
\begin{proof}
\begin{align*}
\int_0^{+\infty}{\frac{\cos \left( ax \right)}{1+x^2}\mathrm{d}x}&=\frac{1}{2}\int_{-\infty}^{+\infty}{\frac{\cos \left( ax \right)}{1+x^2}\mathrm{d}x}=\frac{1}{2}\int_{-\infty}^{+\infty}{\cos \left( ax \right) \left( \int_0^{+\infty}{e^{-\left( 1+x^2 \right) y}\mathrm{d}y} \right) \mathrm{d}x}\\
&=\frac{1}{2}\int_{-\infty}^{+\infty}{\left( \int_0^{+\infty}{e^{-\left( 1+x^2 \right) y}\cos \left( ax \right) \mathrm{d}y} \right) \mathrm{d}x}=\frac{1}{2}\int_{-\infty}^{+\infty}{\int_0^{+\infty}{e^{-\left( 1+x^2 \right) y}\cos \left( ax \right) \mathrm{d}y}}\mathrm{d}x\\
&=\frac{1}{2}\int_0^{+\infty}{\left( \int_{-\infty}^{+\infty}{e^{-\left( 1+x^2 \right) y}\cos \left( ax \right) \mathrm{d}x} \right) \mathrm{d}y}=\frac{1}{2}\int_0^{+\infty}{e^{-y}\left( \int_{-\infty}^{+\infty}{e^{-x^2y}\cos \left( ax \right) \mathrm{d}x} \right)}\mathrm{d}y\\
&=\frac{1}{2}\mathrm{Re}\left( \int_0^{+\infty}{e^{-y}\left( \int_{-\infty}^{+\infty}{e^{-x^2y+\mathrm{i}ax}\mathrm{d}x} \right) \mathrm{d}y} \right) =\frac{1}{2}\mathrm{Re}\left( \int_0^{+\infty}{e^{-y}\left( \int_{-\infty}^{+\infty}{e^{-y\left( x^2-\frac{a\mathrm{i}x}{y}+\left( \frac{a\mathrm{i}}{2y} \right) ^2 \right) -\frac{a^2}{4y}}\mathrm{d}x} \right) \mathrm{d}y} \right) \\
&=\frac{1}{2}\mathrm{Re}\left( \int_0^{+\infty}{e^{-y}\left( \int_{-\infty}^{+\infty}{e^{-y\left( x-\frac{a\mathrm{i}}{2y} \right) ^2-\frac{a^2}{4y}}\mathrm{d}x} \right) \mathrm{d}y} \right) =\frac{1}{2}\mathrm{Re}\left( \int_0^{+\infty}{e^{-y}\left( \int_{-\infty}^{+\infty}{e^{-y\left( x+\frac{a}{2\mathrm{i}y} \right) ^2-\frac{a^2}{4y}}\mathrm{d}x} \right) \mathrm{d}y} \right) \\
&=\frac{1}{2}\mathrm{Re}\left( \int_0^{+\infty}{e^{-y}\left( \int_{-\infty}^{+\infty}{e^{-y\left( x+\frac{a}{2\mathrm{i}y} \right) ^2-\frac{a^2}{4y}}\mathrm{d}\left( x+\frac{a}{2\mathrm{i}y} \right)} \right) \mathrm{d}y} \right) =\frac{1}{2}\mathrm{Re}\left( \int_0^{+\infty}{e^{-y}\left( \int_{-\infty}^{+\infty}{e^{-yx^2-\frac{a^2}{4y}}\mathrm{d}x} \right) \mathrm{d}y} \right) \\
&=\frac{1}{2}\int_0^{+\infty}{e^{-y-\frac{a^2}{4y}}\left( \int_{-\infty}^{+\infty}{e^{-yx^2}\mathrm{d}x} \right) \mathrm{d}y}=\frac{\sqrt{\pi}}{2}\int_0^{+\infty}{\frac{1}{\sqrt{y}}e^{-y-\frac{a^2}{4y}}\mathrm{d}y}\\
&\xlongequal{y=t^2}\sqrt{\pi}\int_0^{+\infty}{e^{-t^2-\frac{a^2}{4t^2}}\mathrm{d}t}\xlongequal{\text{\refexa{example:常用定积分公式}}}\sqrt{\pi}\cdot \frac{\sqrt{\pi}}{2}e^{-|a|}=\frac{\pi}{2}e^{-|a|}.
\end{align*}
\end{proof}

\begin{example}
计算$\int_{0}^{\infty}\frac{1}{(1 + x^{8})^2}\mathrm{d}x$. 
\end{example}
\begin{remark}
由\refpro{proposition:分式化积分-Gamma函数相关性质}可知
对$\forall s>0$,都有
\begin{align*}
\frac{1}{t^s}=\frac{1}{\Gamma \left( s \right)}\int_0^{+\infty}{y^{s-1}e^{-ty}\mathrm{d}y},\forall t\in \mathbb{R}.
\end{align*}
本题的核心想法就是利用上式将$\frac{z}{1+x^8}$转化成积分形式.
\end{remark}
\begin{proof}
注意到
\\
\begin{align*}
\int_0^{+\infty}{ye^{-\left( 1+x^8 \right) y}\mathrm{d}y} \xlongequal{y=\frac{z}{1+x^8}} \frac{1}{(1+x^8)^2}\int_0^{+\infty}{ze^{-z}\mathrm{d}z} = \frac{1}{(1+x^8)^2},
\end{align*}
因此
\\
\begin{align*}
\int_0^{+\infty}{\frac{1}{(1+x^8)^2}\mathrm{d}x} &= \int_0^{+\infty}{\left( \int_0^{+\infty}{ye^{-\left( 1+x^8 \right) y}\mathrm{d}y} \right) \mathrm{d}x} 
= \int_0^{+\infty}{\left( \int_0^{+\infty}{ye^{-\left( 1+x^8 \right) y}\mathrm{d}x} \right) \mathrm{d}y} \\
&= \int_0^{+\infty}{ye^{-y}\left( \int_0^{+\infty}{e^{-x^8y}\mathrm{d}x} \right) \mathrm{d}y} \xlongequal{x=y^{-\frac{1}{8}}z^{\frac{1}{8}}} \int_0^{+\infty}{ye^{-y}\left( \int_0^{+\infty}{y^{-\frac{1}{8}}e^{-z}\mathrm{d}z^{\frac{1}{8}}} \right) \mathrm{d}y} \\
&= \frac{1}{8}\int_0^{+\infty}{y^{\frac{7}{8}}e^{-y}\left( \int_0^{+\infty}{z^{-\frac{7}{8}}e^{-z}\mathrm{d}z} \right) \mathrm{d}y} 
= \frac{1}{8}\int_0^{+\infty}{y^{\frac{7}{8}}e^{-y}\Gamma \left( \frac{1}{8} \right) \mathrm{d}y} \\
&= \frac{1}{8}\Gamma \left( \frac{15}{8} \right) \Gamma \left( \frac{1}{8} \right) = \frac{1}{8} \cdot \frac{7}{8} \Gamma \left( \frac{7}{8} \right) \Gamma \left( \frac{1}{8} \right) \\
&\xlongequal{\hyperref[theorem:余元公式]{\text{余元公式}}} \frac{7\pi}{64 \sin\left( \frac{\pi}{8} \right)} = \frac{7\pi}{32\sqrt{2 - \sqrt{2}}}.
\end{align*}
\end{proof}

\begin{example}
计算积分 $I=\int_{-1}^{2}\frac{1+x^{2}}{1+x^{4}}\,\mathrm{d}x$.
\end{example}
\begin{remark}
在此例中 $I\neq F(2)-F(-1)$. 这是因为 $F$ 并不是 $f$ 在区间 $[-1,2]$ 上的原函数.
\end{remark}
\begin{solution}
在不包含 $0$ 的区间上作变换 $t=x-\frac{1}{x}$ 得
\begin{align*}
\int\frac{1+x^{2}}{1+x^{4}}\,\mathrm{d}x&=\int\frac{x-\frac{1}{x}}{2+\left(x-\frac{1}{x}\right)^{2}}\,\mathrm{d}x=\int\frac{\mathrm{d}t}{2+t^{2}}\\
&=\frac{1}{\sqrt{2}}\arctan\frac{t}{\sqrt{2}}+C=\frac{1}{\sqrt{2}}\arctan\frac{x^{2}-1}{\sqrt{2}x}+C.
\end{align*}
这说明在区间 $[-1,0)$ 和 $(0,2]$ 上, 函数 $f(x)=\frac{1+x^{2}}{1+x^{4}}$ 的一个原函数是
$$F(x)=\frac{1}{\sqrt{2}}\arctan\frac{x^{2}-1}{\sqrt{2}x}.$$
因此
\begin{align*}
\int_{-1}^{0}f(x)\,\mathrm{d}x&=F(0^{-})-F(-1)=\frac{\pi}{2\sqrt{2}}-0=\frac{\pi}{2\sqrt{2}},\\
\int_{0}^{2}f(x)\,\mathrm{d}x&=F(2)-F(0^{+})=\frac{1}{\sqrt{2}}\arctan\frac{3}{2\sqrt{2}}+\frac{\pi}{2\sqrt{2}}.
\end{align*}
故
$$I=\int_{-1}^{0}f(x)\,\mathrm{d}x+\int_{0}^{2}f(x)\,\mathrm{d}x=\frac{\pi}{\sqrt{2}}+\frac{1}{\sqrt{2}}\arctan\frac{3}{2\sqrt{2}}.$$

\end{solution}

\begin{example}
计算$I\left( b \right) =\int_0^{\frac{\pi}{2}}{\frac{1}{b+\tan x}\mathrm{d}x}.$
\end{example}
\begin{solution}

{\color{blue} 解法一:}注意到
\begin{align*}
I\left( b \right) =\int_0^{\frac{\pi}{2}}{\frac{1}{b+\frac{\sin x}{\cos x}}\mathrm{d}x}=\int_0^{\frac{\pi}{2}}{\frac{\cos x}{\sin x+b\cos x}\mathrm{d}x},
\end{align*}
记
\begin{align*}
I=\int_0^{\frac{\pi}{2}}{\frac{\cos x}{\sin x+b\cos x}\mathrm{d}x},\quad J=\int_0^{\frac{\pi}{2}}{\frac{\sin x}{\sin x+b\cos x}\mathrm{d}x}.
\end{align*}
从而
\begin{align*}
bI+J=\int_0^{\frac{\pi}{2}}{\frac{\sin x+b\cos x}{\sin x+b\cos x}\mathrm{d}x}=\frac{\pi}{2},
\end{align*}
\begin{align*}
I-bJ=\int_0^{\frac{\pi}{2}}{\frac{\cos x-b\sin x}{\sin x+b\cos x}\mathrm{d}x}=\ln \left( \sin x+b\cos x \right) \big| _{0}^{\frac{\pi}{2}}=-\ln b.
\end{align*}
联立上面两式解得
\begin{align*}
I\left( b \right) =I=\frac{\pi b}{2\left( b^2+1 \right)}-\frac{\ln b}{b^2+1}.
\end{align*}

{\color{blue} 解法二:}令$t=\tan x,$则
\begin{align*}
I\left( b \right) =\int_0^{\frac{\pi}{2}}{\frac{1}{b+\tan x}\mathrm{d}x}=\int_0^{+\infty}{\frac{1}{\left( b+t \right) \left( 1+t^2 \right)}\mathrm{d}t}=\lim\limits_{A\rightarrow +\infty}\int_0^A{\frac{1}{\left( b+t \right) \left( 1+t^2 \right)}\mathrm{d}t}.
\end{align*}
注意到
\begin{align*}
\int_0^A{\frac{1}{\left( b+t \right) \left( 1+t^2 \right)}\mathrm{d}t}&=\frac{1}{b^2+1}\int_0^A{\frac{1}{b+t}\mathrm{d}t}+\frac{1}{b^2+1}\int_0^A{\frac{-t+b}{1+t^2}\mathrm{d}t} \\
&=\frac{1}{b^2+1}\ln \frac{b+A}{b}-\frac{1}{2\left( b^2+1 \right)}\int_0^A{\frac{2t-2b}{1+t^2}\mathrm{d}t} \\
&=\frac{1}{b^2+1}\ln \frac{b+A}{b}-\frac{1}{2\left( b^2+1 \right)}\int_0^A{\frac{2t}{1+t^2}\mathrm{d}t}+\frac{b}{b^2+1}\int_0^A{\frac{1}{1+t^2}\mathrm{d}t} \\
&=\frac{1}{b^2+1}\ln \frac{b+A}{b}-\frac{1}{2\left( b^2+1 \right)}\ln \left( 1+A^2 \right) +\frac{b}{b^2+1}\arctan A \\
&=\frac{1}{b^2+1}\ln \frac{b+A}{b\sqrt{1+A^2}}+\frac{b}{b^2+1}\arctan A,
\end{align*}
于是
\begin{align*}
I\left( b \right) &=\lim\limits_{A\rightarrow +\infty}\int_0^A{\frac{1}{\left( b+t \right) \left( 1+t^2 \right)}\mathrm{d}t} \\
&=\lim\limits_{A\rightarrow +\infty}\left[ \frac{1}{b^2+1}\ln \frac{b+A}{b\sqrt{1+A^2}}+\frac{b}{b^2+1}\arctan A \right] \\
&=\frac{\pi b}{2\left( b^2+1 \right)}-\frac{\ln b}{b^2+1}.
\end{align*}
\end{solution}

\begin{example}
设 $\alpha\in(1,2)$，证明
\begin{align*}
\int_{\mathbb{R}}\left(\frac{2}{\alpha(1+x^2)^{1+\frac{1}{\alpha}}}-\frac{1}{(1+x^2)^{\frac{1}{\alpha}}}\right)\mathrm{d}x<0.
\end{align*}
\end{example}
\begin{proof}
令 $x=\tan\theta$，则 $\mathrm{d}x=\sec^2\theta\,\mathrm{d}\theta$，且 $\theta\in(-\frac{\pi}{2},\frac{\pi}{2})$，$1+x^2=\sec^2\theta$。代入得
\begin{align*}
I=\int_{-\frac{\pi}{2}}^{\frac{\pi}{2}}\left[\frac{2}{\alpha}\sec^{-2(1+\frac{1}{\alpha})}\theta\cdot\sec^2\theta-\sec^{-\frac{2}{\alpha}}\theta\cdot\sec^2\theta\right]\mathrm{d}\theta 
=\int_{-\frac{\pi}{2}}^{\frac{\pi}{2}}\left[\frac{2}{\alpha}\cos^{\frac{2}{\alpha}}\theta-\cos^{\frac{2}{\alpha}-2}\theta\right]\mathrm{d}\theta.
\end{align*}
被积函数是偶函数，故
\begin{align*}
I=2\int_0^{\frac{\pi}{2}}\left[\frac{2}{\alpha}\cos^{\frac{2}{\alpha}}\theta-\cos^{\frac{2}{\alpha}-2}\theta\right]\mathrm{d}\theta.
\end{align*}
记 $p=\frac{2}{\alpha}$。因为 $\alpha\in(1,2)$，所以 $p\in(1,2)$。于是
\begin{align*}
I=2\int_0^{\frac{\pi}{2}}\left(p\cos^p\theta-\cos^{p-2}\theta\right)\mathrm{d}\theta.
\end{align*}
令
\begin{align*}
A_m\triangleq \int_0^{\frac{\pi}{2}}\cos^m\theta\,\mathrm{d}\theta,\quad m>-1.
\end{align*}
则利用分部积分可得,对$\forall m>-1$,有
\begin{align*}
A_m&=\int_0^{\frac{\pi}{2}}{\cos ^m\theta \,\mathrm{d}\theta}=\sin \theta \cos ^{m-1}\theta \Big |_{0}^{\frac{\pi}{2}}+\left( m-1 \right) \int_0^{\frac{\pi}{2}}{\sin ^2\theta \cos ^{m-2}\theta \,\mathrm{d}\theta}
\\
&=\left( m-1 \right) \int_0^{\frac{\pi}{2}}{\left( 1-\cos ^2\theta \right) \cos ^{m-2}\theta \,\mathrm{d}\theta}=\left( m-1 \right) A_{m-2}-\left( m-1 \right) A_m.
\end{align*}
进而
\begin{align*}
A_m=\frac{m-1}{m}A_{m-2}\quad \left( m>1 \right) .
\end{align*}
将 $m=p$ 代入得
\begin{align*}
A_p=\frac{p-1}{p}A_{p-2}\Longrightarrow pA_p=(p-1)A_{p-2}.
\end{align*}
因此
\begin{align*}
I=2\left(pA_p-A_{p-2}\right)
=2\left((p-1)A_{p-2}-A_{p-2}\right) =2(p-2)A_{p-2}.
\end{align*}
因为 $p\in(1,2)$，故 $p-2<0$。而
\begin{align*}
A_{p-2}=\int_0^{\frac{\pi}{2}}\cos^{p-2}\theta\,\mathrm{d}\theta>0
\end{align*}
所以
\begin{align*}
I=2(p-2)A_{p-2}<0.
\end{align*}
\end{proof}

\begin{proposition}\label{proposition:常见积分结论111}
证明:
\begin{align*}
\displaystyle\int_0^{+\infty} e^{-x}\sin(xt)\,\mathrm{d}x = \frac{t}{1+t^2}.
\end{align*}
\end{proposition}
\begin{proof}
令 $I = \int_0^{+\infty} e^{-x}\sin(xt)\,\mathrm{d}x$.
分部积分，得
\begin{align*}
I = \left[ -e^{-x}\sin(xt) \right]_0^{+\infty} + t\int_0^{+\infty} e^{-x}\cos(xt)\,\mathrm{d}x = t\int_0^{+\infty} e^{-x}\cos(xt)\,\mathrm{d}x = t\,J,
\end{align*}
其中 $J = \int_0^{+\infty} e^{-x}\cos(xt)\,\mathrm{d}x$.
对 $J$ 再进行分部积分得
\begin{align*}
J = \left[ -e^{-x}\cos(xt) \right]_0^{+\infty} - t\int_0^{+\infty} e^{-x}\sin(xt)\,\mathrm{d}x = 1 - t I.
\end{align*}
将 $J = 1 - t I$ 代入 $I = t J$ 得
\begin{align*}
I = t(1 - t I) = t - t^2 I,
\end{align*}
移项整理得 $I(1 + t^2) = t$, 故
\begin{align*}
I = \frac{t}{1+t^2}.
\end{align*}
\end{proof}

\begin{example}
求反常积分
\begin{align*}
\int_0^{+\infty} e^{-x}\cdot \frac{1-\cos x}{x}\,\mathrm{d}x
\end{align*}
的值.
\end{example}
\begin{solution}

{\color{blue}解法一:} 化为二重积分.
利用恒等式
\begin{align*}
\frac{1-\cos x}{x} = \int_0^1 \sin(xt)\,\mathrm{d}t,
\end{align*}
代入原积分并交换积分次序得
\begin{align*}
\int_0^{+\infty} e^{-x}\cdot \frac{1-\cos x}{x}\,\mathrm{d}x
= \int_0^{+\infty} e^{-x} \biggl( \int_0^1 \sin(xt)\,\mathrm{d}t \biggr) \mathrm{d}x
= \int_0^1 \biggl( \int_0^{+\infty} e^{-x}\sin(xt)\,\mathrm{d}x \biggr) \mathrm{d}t.
\end{align*}
由\cref{proposition:常见积分结论111}知
\begin{align*}
\int_0^{+\infty} e^{-x}\sin(xt)\,\mathrm{d}x = \frac{t}{1+t^2},
\end{align*}
所以
\begin{align*}
\int_0^{+\infty} e^{-x}\cdot \frac{1-\cos x}{x}\,\mathrm{d}x
= \int_0^1 \frac{t}{1+t^2}\,\mathrm{d}t
= \frac{1}{2}\ln(1+t^2)\Big|_0^1
= \frac{1}{2}\ln 2.
\end{align*}

{\color{blue}解法二:}
令
\begin{align*}
J(b) = \int_0^{+\infty} e^{-x}\cdot \frac{1-\cos(bx)}{x}\,\mathrm{d}x, \quad J(0)=0.
\end{align*}
对 $b$ 求导，再利用\cref{proposition:常见积分结论111}得
\begin{align*}
J'(b) = \int_0^{+\infty} e^{-x}\sin(bx)\,\mathrm{d}x = \frac{b}{1+b^2}.
\end{align*}
积分得
\begin{align*}
J(b) = \int_0^b \frac{t}{1+t^2}\,\mathrm{d}t = \frac{1}{2}\ln(1+b^2).
\end{align*}
原积分即为 $b=1$ 时的值:
\begin{align*}
\int_0^{+\infty} e^{-x}\cdot \frac{1-\cos x}{x}\,\mathrm{d}x = J(1) = \frac{1}{2}\ln 2.
\end{align*}
\end{solution}





\end{document}